\section{事故機構・設計運用と個別責任の区別}

\subsection{平成30年度 問14}
\label{q:h30-reg-apt-14}

\noindent\textbf{原典：} sources/H30/aptitude.pdf、PDF物理14頁。

\textbf{公式正答：} ALL（all\_awarded）。

\par\textbf{独立解答と公式正答に相違あり：} 独立解答 2|5。初回2|5に対し公式ALL。2の誤った有罪は最高裁2017年決定で照合、5の不適格は原典公式注記で確認。1・3・4の個別判決・量刑比較は一次裁判資料未確認。全肢真とせず部分校正、Astra難問査読へ。

\paragraph{問題の要旨} 5事故の刑事責任を判決と比較する設問と、笹子の不起訴を判決扱いした不適切出題注記を検討する。

\Needspace{4\baselineskip}\paragraph{確認した条件・記述（独自要約）}\mbox{}\par

\noindent\textbf{方向：} 実際の刑事判決と異なる事例を選ぶ。\par

\noindent\textbf{原典注記：} 選択肢5は不起訴処分で判決ではなく、試験日現在検察審査会への申立てがあり不適切出題。\par

\noindent\textbf{時点：} 平成30年の試験日。後年の経過から逆算しない。\par

\noindent\textbf{OCR補正：} 選択肢1はシンドラー社元社員。\par

\paragraph{問われる核心} 事故の技術的原因と個別の刑事責任・採点措置の区別

\paragraph{必要知識} 事故の技術的原因と個別の刑事責任・採点措置の区別。初回独立候補2|5を保持する。2の有罪は、歴代3社長を対象とした最高裁の2017年決定と異なる。5は、公式問題の注記が示すとおり不起訴処分を刑事裁判の判決と同列に置いている。公式正答は不適切出題による全員正解（ALL）であり、初回候補をALLへ上書きしない。1・3・4の個別判決の一次資料は未確認で、全肢を確定済みにしない。全員正解は採点上の措置であり、5事例の記述が全て正しいという意味ではない。

\paragraph{考え方・正答を選ぶ理由（査読待ち初稿）} 初回独立候補2|5を保持する。2の有罪は、歴代3社長を対象とした最高裁の2017年決定と異なる。5は、公式問題の注記が示すとおり不起訴処分を刑事裁判の判決と同列に置いている。公式正答は不適切出題による全員正解（ALL）であり、初回候補をALLへ上書きしない。1・3・4の個別判決の一次資料は未確認で、全肢を確定済みにしない。全員正解は採点上の措置であり、5事例の記述が全て正しいという意味ではない。

\Needspace{4\baselineskip}\paragraph{選択肢の検討（査読待ち初稿）}\mbox{}\par

\noindent\textbf{1：判決一次資料要確認} 2006年エレベーター事故で、シンドラー社元社員の刑事責任はないとされた。。原画像の主語は元社員であり、OCRの社員を補正する。個別被告人の判決と確定段階を示す一次の裁判資料は未確認であり、会社・全保守関係者の責任へ一般化しない。\par

\noindent\textbf{2：判決と異なる} 2005年福知山線事故で、歴代3社長が有罪となった。。2017年6月12日最高裁第二小法廷決定、平成27年（あ）741号は、歴代3社長への無罪の第一審を是認した原判断を相当とし、各上告を棄却した。有罪という結末が異なる。\par

\noindent\textbf{3：判決一次資料要確認} 2004年回転ドア事故で、メーカー営業開発部長の刑事責任が森ビル関係者より重いとされた。。事故防止上の製造者・管理者の役割と、個別被告人への量刑比較は別問題である。刑の重さと理由を一次の判決文で照合できておらず、適切と確定しない。\par

\noindent\textbf{4：判決一次資料要確認} 2000年低脂肪乳食中毒事故で、工場関係者が有罪となった。。厚労省資料は停電時の菌増殖・エンテロトキシンA産生と脱脂粉乳の因果を示すが、工場関係者の有罪判決そのものではない。刑事事件の一次記録は追加確認を要する。\par

\noindent\textbf{5：公式注記による不適格} 2012年笹子トンネル事故で、社長らは予見できないとして刑事責任なしとする判決を受けた。。原典に付された公式説明は、不起訴処分で刑事裁判に当たらず、試験日現在は検察審査会への審査申立てもあったため不適格とする。不起訴を無罪判決に読み替えられない。\par

\paragraph{再利用する考え方} 対象、条件、時点を先に確認する。事故の技術的原因と個別の刑事責任・採点措置の区別を具体的な事例へ適用し、技術的原因・法的要件・倫理的判断のどれを問うか区別する。

\paragraph{共通理論} \hyperref[theory:safety:responsibility]{第\ref*{theory:safety:responsibility}節}, \hyperref[theory:safety:cases]{第\ref*{theory:safety:cases}節}

\par 参考：\cite{court_jr2017,jtsb_fukuchiyama2007,mhlw_milk2002,mlit_sasago2013}。

\paragraph{関連問題} \hyperref[q:r05-reg-apt-09]{R05-REG-APT-09}

\par 難易度 3（校準初稿）、優先度未評価。

\subsection{令和5年度 問9}
\label{q:r05-reg-apt-09}

\noindent\textbf{原典：} sources/R05/aptitude.pdf、PDF物理9頁。

\textbf{公式正答：} 4（single）。

\paragraph{問題の要旨} 食中毒、回転ドア、鉄道、エレベーター、トンネル天井板の事故原因を比較する。

\Needspace{4\baselineskip}\paragraph{確認した条件・記述（独自要約）}\mbox{}\par

\noindent\textbf{方向：} 事実と異なるもの。\par

\noindent\textbf{事故年：} 2000、2004、2005、2006、2012。\par

\noindent\textbf{条件：} 鉄道のATSは当該急曲線の速度超過防止。エレベーターの運動方向を実画像と政府事故調査で確認。\par

\paragraph{問われる核心} 事故機構の再現と設計・運用・組織の因果の区別

\paragraph{必要知識} 事故機構の再現と設計・運用・組織の因果の区別。4は事故の機構を逆にしている。2006年の事象は、かごが戸開状態で不意に上昇した際の挟まれであり、ケーブル破断・落下ではない。1〜3・5は、直接の物理的事象だけでなく、その前提を維持できなかった設計・運用・組織上の条件も述べる。鉄道のATSの機能と、組織背景の推定の強さを区別して4を選ぶ。

\paragraph{考え方・正答を選ぶ理由（査読待ち初稿）} 4は事故の機構を逆にしている。2006年の事象は、かごが戸開状態で不意に上昇した際の挟まれであり、ケーブル破断・落下ではない。1〜3・5は、直接の物理的事象だけでなく、その前提を維持できなかった設計・運用・組織上の条件も述べる。鉄道のATSの機能と、組織背景の推定の強さを区別して4を選ぶ。

\Needspace{4\baselineskip}\paragraph{選択肢の検討（査読待ち初稿）}\mbox{}\par

\noindent\textbf{1：適切（事故対応の説明）} 2000年食中毒は停電後の不適切な処置と情報開示・回収の遅れ等が関係した。。厚労省は停電時の菌増殖・毒素産生を説明し、当事企業も汚染原料と対応の問題を振り返る。個別の経営者の刑事責任を認定する記述ではない。\par

\noindent\textbf{2：適切（安全の教訓）} 2004年回転ドアは検知・停止の設計と運用、軽量設計の軽視等に問題があった。。検知・非常停止だけに安全を委ね、重量・慣性と使用条件を軽視すると停止前の挟まれ危険が残る。大学の技術者倫理事例集は軽量設計からの変更を扱う。事故の一般的教訓と個別被告人の判決を分ける。\par

\noindent\textbf{3：適切（機能・推定の限定に注意）} 2005年福知山線は曲線での速度超過・ブレーキの遅れと安全確保の組織的問題が関係した。。事故調査報告の直接原因は制限速度を大幅に超えた進入で、運転士管理方法の関与は可能性として記す。原典のATS未設置は当該曲線の速度超過防止機能を指し、列車の全ATS不搭載を意味しない。\par

\noindent\textbf{4：事実と異なる} 2006年エレベーター事故はケーブルが破断してかごが落下した。。公式の事故調査は、戸が開いた状態でかごが上昇し利用者が挟まれた事象を扱う。ケーブル破断による落下とした故障機構・運動方向が事実と異なる。\par

\noindent\textbf{5：適切（複合的要因）} 2012年笹子は天井板の接着系アンカー等の脱落、設計・施工・維持管理が関係した。。国交省委員会は接着系アンカーの引抜き等と、設計・施工・点検の課題を整理する。単純な材料の経年劣化だけに原因を限定しない。\par

\paragraph{再利用する考え方} 対象、条件、時点を先に確認する。事故機構の再現と設計・運用・組織の因果の区別を具体的な事例へ適用し、技術的原因・法的要件・倫理的判断のどれを問うか区別する。

\paragraph{共通理論} \hyperref[theory:safety:cases]{第\ref*{theory:safety:cases}節}, \hyperref[theory:safety:barriers]{第\ref*{theory:safety:barriers}節}

\par 参考：\cite{mhlw_milk2002,milk_incident_company,muroran_ethics_cases,jtsb_fukuchiyama2007,caa_elevator2006,mlit_sasago2013}。

\paragraph{関連問題} \hyperref[q:h30-reg-apt-14]{H30-REG-APT-14}

\par 難易度 2（校準初稿）、優先度未評価。

\section{ALARPの水準・実行可能性・著しい不均衡}

\subsection{令和5年度 問14}
\label{q:r05-reg-apt-14}

\noindent\textbf{原典：} sources/R05/aptitude.pdf、PDF物理14頁。

\textbf{公式正答：} 3（single）。

\par\textbf{検討未確定：} 肢2「あらゆるリスク」の範囲とGuide51のALARP水準の条項の直接照合が未完。英国の法的ALARP適用と3領域図の説明も同一に扱わない。公式3・初回3は保持。

\paragraph{問題の要旨} ALARPを事業予算の上限に置換せず、本質安全設計と警告の役割を整理する。

\Needspace{4\baselineskip}\paragraph{確認した条件・記述（独自要約）}\mbox{}\par

\noindent\textbf{判断方向：} 各肢の理由に記録。組合せ・個数問題では記述の真偽と回答番号を区別する。\par

\noindent\textbf{原典照合：} 全5候補、表の列順、否定、条件、規格番号、法令時点を23原典画像で照合。\par

\paragraph{問われる核心} ALARPを事業予算の上限に置換せず、本質安全設計と警告の役割を整理する。

\paragraph{必要知識} 追加低減の可能性、費用対効果の著しい不均衡、予算との差、警告。

\paragraph{考え方・正答を選ぶ理由（査読待ち初稿）} 最も不適切な3を選ぶ。ALARPはどの追加低減が合理的に実行可能かを評価する原理であり、経営者が設定した製品原価・予算だけで決めるものではない。2の全称は適用水準と法的枠組みで異なり、原規格を確認せず一律の真偽を補わない。

\Needspace{4\baselineskip}\paragraph{選択肢の検討（査読待ち初稿）}\mbox{}\par

\noindent\textbf{1：適切（条件内）} 反復と残留リスクの妥当性・文書。。対策後にも評価し、判断根拠を残す。\par

\noindent\textbf{2：適切（条件内）} あらゆるリスクを合理的可能性まで低減。。適用枠組みの低減要求としては理解できるが、「あらゆる」と水準の限定の関係は条文を確認する必要がある。\par

\noindent\textbf{3：不適切} 低減が不能又は当初予算に収まらなければ許容。。予算枠は組織の計画で、効果に対する著しい費用・負担の不均衡とは別である。僅かに予算を超えるだけでは許容の根拠にならない。\par

\noindent\textbf{4：適切（条件内）} ステップ1は本質的設計で最重要。。危険源自体を減らすことを先に検討する。情報はこの検討の代替ではない。\par

\noindent\textbf{5：適切（条件内）} 警告は明瞭、可読、消えにくく理解可能。。読めて理解され、使用期間に機能する必要がある。表示が存在するだけで十分ではない。\par

\paragraph{再利用する考え方} 危険源、曝露、損害、評価基準、対策、残留リスクを分ける。全称・必須・予算だけ・優先順位なしという限定を検討し、出題当時と現在の規定を混同しない。

\paragraph{共通理論} \hyperref[theory:safety:alarp]{第\ref*{theory:safety:alarp}節}, \hyperref[theory:safety:hierarchy]{第\ref*{theory:safety:hierarchy}節}, \hyperref[theory:safety:warnings]{第\ref*{theory:safety:warnings}節}

\par 参考：\cite{saf_ra2006,saf_parliament2006,saf_guide51}。

\paragraph{関連問題} \hyperref[q:r03-reg-apt-10]{R03-REG-APT-10}

\par 難易度 2（校準初稿）、優先度未評価。

\subsection{令和7年度 問14}
\label{q:r07-reg-apt-14}

\noindent\textbf{原典：} sources/R07/aptitude.pdf、PDF物理14頁。

\textbf{公式正答：} 2（single）。

\par\textbf{検討未確定：} 前文の安全＝広く受け入れ可能なレベルという等置は、中間の条件付き許容を含む安全概念との整合性に疑義。図と候補は確認し公式2・初回2を保持。

\paragraph{問題の要旨} ALARP領域の条件付き許容を、広く受容される低リスクや事業予算と区別する。

\Needspace{4\baselineskip}\paragraph{確認した条件・記述（独自要約）}\mbox{}\par

\noindent\textbf{判断方向：} 各肢の理由に記録。組合せ・個数問題では記述の真偽と回答番号を区別する。\par

\noindent\textbf{原典照合：} 全5候補、表の列順、否定、条件、規格番号、法令時点を23原典画像で照合。\par

\paragraph{問われる核心} ALARP領域の条件付き許容を、広く受容される低リスクや事業予算と区別する。

\paragraph{必要知識} 3領域、条件付き許容、著しい不均衡、損害と補償、予算。

\paragraph{考え方・正答を選ぶ理由（査読待ち初稿）} 最も適切な2を選ぶ。上の許容不能、中間の条件付許容、下の広く受容という3水準を分ける。中間では追加低減の合理性を評価する。前文が「受容できないリスクなし」を下の水準だけと等置する点には、中間の条件付き許容との整合性の疑義を残す。

\Needspace{4\baselineskip}\paragraph{選択肢の検討（査読待ち初稿）}\mbox{}\par

\noindent\textbf{1：不適切} 領域内は無条件で許容。。領域内にあるだけでは足りず、合理的に可能な低減を尽くす条件を伴う。\par

\noindent\textbf{2：条件に適合} 追加低減の費用が改善と全く釣り合わない。。更なる低減が実行不能又は費用・負担が効果に対し著しく不均衡という条件である。単なる少額の超過ではない。\par

\noindent\textbf{3：不適切} 賠償費用より対策費が高い。。補償金の支出だけを便益とせず、危害回避・生命健康・時間等の負担を検討する。\par

\noindent\textbf{4：領域が違う} 残存リスクが無視できる。。無視できる低水準は図の下の広く受け入れられる領域に対応し、問う中間水準の条件とは違う。\par

\noindent\textbf{5：不適切} 予算を別事業に使いたい。。他の事業へ資金を使いたいという選好だけでは、合理的に可能な安全措置を省略する根拠にならない。\par

\paragraph{再利用する考え方} 危険源、曝露、損害、評価基準、対策、残留リスクを分ける。全称・必須・予算だけ・優先順位なしという限定を検討し、出題当時と現在の規定を混同しない。

\paragraph{共通理論} \hyperref[theory:safety:alarp]{第\ref*{theory:safety:alarp}節}, \hyperref[theory:safety:acceptability]{第\ref*{theory:safety:acceptability}節}

\par 参考：\cite{saf_parliament2006,saf_ra2006}。

\paragraph{関連問題} \hyperref[q:r04-reg-apt-06]{R04-REG-APT-06}

\par 難易度 2（校準初稿）、優先度未評価。

\section{最適化・意思決定}

\subsection{令和4年度 問7}
\label{q:r04-reg-apt-07}

\noindent\textbf{原典：} sources/R04/aptitude.pdf、PDF物理6頁。

\textbf{公式正答：} 3（single）。

\par\textit{原典本文の校正と独自解説は未作成。分類のみの登録項目。}

\section{材料の応力・変形}

\subsection{令和元年度 問13}
\label{q:r01-reg-apt-13}

\noindent\textbf{原典：} sources/R01/aptitude.pdf、PDF物理13頁。

\textbf{公式正答：} 4（single）。

\par\textit{原典本文の校正と独自解説は未作成。分類のみの登録項目。}

\section{公益・説明責任}

\subsection{平成27年度 問4}
\label{q:h27-reg-apt-04}

\noindent\textbf{原典：} sources/H27/aptitude.pdf、PDF物理4頁。

\textbf{公式正答：} 3（single）。

\par\textit{原典本文の校正と独自解説は未作成。分類のみの登録項目。}

\subsection{平成27年度 問5}
\label{q:h27-reg-apt-05}

\noindent\textbf{原典：} sources/H27/aptitude.pdf、PDF物理5頁。

\textbf{公式正答：} 5（single）。

\par\textit{原典本文の校正と独自解説は未作成。分類のみの登録項目。}

\subsection{平成28年度 問3}
\label{q:h28-reg-apt-03}

\noindent\textbf{原典：} sources/H28/aptitude.pdf、PDF物理3頁。

\textbf{公式正答：} 1（single）。

\par\textit{原典本文の校正と独自解説は未作成。分類のみの登録項目。}

\subsection{平成28年度 問9}
\label{q:h28-reg-apt-09}

\noindent\textbf{原典：} sources/H28/aptitude.pdf、PDF物理9頁。

\textbf{公式正答：} 5（single）。

\par\textit{原典本文の校正と独自解説は未作成。分類のみの登録項目。}

\subsection{令和3年度 問2}
\label{q:r03-reg-apt-02}

\noindent\textbf{原典：} sources/R03/aptitude.pdf、PDF物理2頁。

\textbf{公式正答：} 4（single）。

\par\textit{原典本文の校正と独自解説は未作成。分類のみの登録項目。}

\section{安全・リスク評価}

\subsection{平成27年度 問14}
\label{q:h27-reg-apt-14}

\noindent\textbf{原典：} sources/H27/aptitude.pdf、PDF物理14頁。

\textbf{公式正答：} 5（single）。

\par\textit{原典本文の校正と独自解説は未作成。分類のみの登録項目。}

\subsection{令和2年度 問9}
\label{q:r02-reg-apt-09}

\noindent\textbf{原典：} sources/R02/aptitude.pdf、PDF物理9頁。

\textbf{公式正答：} 3（single）。

\par\textit{原典本文の校正と独自解説は未作成。分類のみの登録項目。}

\section{災害リスク情報・避難制度と事前準備}

\subsection{令和元年度再試験 問13}
\label{q:r01-ret-apt-13}

\noindent\textbf{原典：} sources/R01\_retest/aptitude.pdf、PDF物理14頁。

\textbf{公式正答：} 1（single）。

\paragraph{問題の要旨} 水害・土砂災害で地図、警報、避難情報、事前準備、自主判断の5記述を検討する。

\Needspace{4\baselineskip}\paragraph{確認した条件・記述（独自要約）}\mbox{}\par

\noindent\textbf{方向：} 不適切な記述の数。\par

\noindent\textbf{時点：} 2019年度再試験。旧避難情報名称を現在へ書き換えない。\par

\noindent\textbf{導入部：} AR5の評価した排出シナリオ、1901〜2017日本の降水日数の記述は該当一次資料で確認。すべての仮想シナリオや全地点個別の結果へ拡張しない。\par

\noindent\textbf{ア：} ハザードマップで地域の災害範囲・程度・避難情報を知る。 ○：対象災害や想定条件、経路と場所を事前に理解する。白い範囲が全災害で安全という意味ではない。\par

\noindent\textbf{イ：} 数日前の情報から注意報・警報・特別警報へ危険に応じ発表。 ○：当時の情報体系の概要として適切。ただし急変で常にすべての段階を順次通る保証はない。\par

\noindent\textbf{ウ：} 市町村長の準備・高齢者等避難開始、勧告、指示等で避難。 ○：2019年度の名称として適切。勧告は2021年5月20日に廃止し指示へ一本化。\par

\noindent\textbf{エ：} 備蓄・持出品、安否確認・集合、保険等を事前に準備。 ○：避難、家族との連絡、被害後の生活を異なる対策で支える。\par

\noindent\textbf{オ：} 発令が間に合わなくても危険なら自ら避難判断。 ○：行政情報の受信だけに依存すると急変に間に合わない。自ら得られる危険情報も判断の入力とする。\par

\paragraph{問われる核心} 災害情報を行動に接続する事前準備と制度時点

\paragraph{必要知識} 災害情報を行動に接続する事前準備と制度時点。ア〜オはすべて適切で不適切は0件、候補1。原典の当時の情報名を確認した上で、地図の想定外・突然の危険・警報発表主体を補足する。現在は2021年の避難勧告廃止と、2026年5月29日の警戒レベルを付した新しい防災気象情報を区別して学ぶ。

\paragraph{考え方・正答を選ぶ理由（査読待ち初稿）} ア〜オはすべて適切で不適切は0件、候補1。原典の当時の情報名を確認した上で、地図の想定外・突然の危険・警報発表主体を補足する。現在は2021年の避難勧告廃止と、2026年5月29日の警戒レベルを付した新しい防災気象情報を区別して学ぶ。

\Needspace{4\baselineskip}\paragraph{選択肢の検討（査読待ち初稿）}\mbox{}\par

\noindent\textbf{1：整合する} 0。全5記述が当時の一般的備えとして適切で、不適切0と一致。\par

\noindent\textbf{2：整合しない} 1。不適切な記述を1件とする根拠がない。\par

\noindent\textbf{3：整合しない} 2。不適切な記述を2件とする根拠がない。\par

\noindent\textbf{4：整合しない} 3。不適切な記述を3件とする根拠がない。\par

\noindent\textbf{5：整合しない} 4。不適切な記述を4件とする根拠がない。\par

\paragraph{再利用する考え方} 活動・物質・制度の対象と評価の境界を固定し、用語の入替、全部/一部、法的義務/目標、出題時/現在を区別する。災害情報を行動に接続する事前準備と制度時点を個別の技術活動に結び付ける。

\paragraph{共通理論} \hyperref[theory:safety:preparedness]{第\ref*{theory:safety:preparedness}節}

\par 参考：\cite{env_ar5,env_rain2017,env_evac2021,env_warning2026}。

\par 難易度 2（校準初稿）、優先度未評価。

\section{インフラの診断記録と協働}

\subsection{平成26年度 問13}
\label{q:h26-reg-apt-13}

\noindent\textbf{原典：} sources/H26/aptitude.pdf、PDF物理13頁。

\textbf{公式正答：} 5（single）。

\par\textbf{検討未確定：} H26時点の全分野・地域の制度を網羅した確認ではない。協働例で全禁止を否定するが、専門点検や危険作業を無条件に住民へ移せるとはしない。

\paragraph{問題の要旨} 維持管理の行政責任と住民・民間の協働を区別する。

\Needspace{4\baselineskip}\paragraph{確認した条件・記述（独自要約）}\mbox{}\par

\noindent\textbf{原典照合：} 22原典画像で25問の全五肢・数値・単位・否定を確認。\par

\noindent\textbf{対象：} 維持管理の行政責任と住民・民間の協働を区別する。\par

\paragraph{問われる核心} 維持管理の行政責任と住民・民間の協働を区別する。

\paragraph{必要知識} 5は行政の管理責任から地域住民・民間の参加を全面禁止へ飛躍した点が不適切。施設の安全を管理する専門的責任を明確にしつつ、点検情報・清掃等で役割を分担できる。維持管理は点検・診断・措置・記録を次の周期へつなぐ活動。

\paragraph{考え方・正答を選ぶ理由（査読待ち初稿）} 5は行政の管理責任から地域住民・民間の参加を全面禁止へ飛躍した点が不適切。施設の安全を管理する専門的責任を明確にしつつ、点検情報・清掃等で役割を分担できる。維持管理は点検・診断・措置・記録を次の周期へつなぐ活動。

\Needspace{4\baselineskip}\paragraph{選択肢の検討（査読待ち初稿）}\mbox{}\par

\noindent\textbf{1：代表的条件で適切} 高度成長期の整備と老朽化費増大。整備時期の集中が将来の点検・修繕・更新の需要を増やす。年数だけで個別施設の健全度を決めない。\par

\noindent\textbf{2：代表的条件で適切} 計画で費用平準化・性能保持。予防・計画的な修繕でライフサイクル全体の需要と性能を調整する。必ず費用が一定になる保証ではない。\par

\noindent\textbf{3：代表的条件で適切} 人材確保・技術向上・新技術。点検・診断・施工の能力と技術を確保する。ロボット等を導入しても責任と評価を省略しない。\par

\noindent\textbf{4：代表的条件で適切} 点検履歴・施設諸元の情報整備。情報を記録して次の診断・措置の根拠とする。データ未整備の把握も必要。\par

\noindent\textbf{5：不適切} 住民自主管理・民間ボランティアは禁止。管理者の責任と市民参加は両立する。国交省には道路のアダプト制度等の協働例があり、全禁止ではない。\par

\paragraph{再利用する考え方} 点検・診断・措置・記録のサイクルを再構成し、施設ごとの基準と全施設一律の主張を区別する。管理者の専門的責任と住民との協働を両立させる。

\paragraph{共通理論} \hyperref[theory:safety:maintenance]{第\ref*{theory:safety:maintenance}節}

\par 参考：\cite{eng52_infra2014,eng52_adopt}。

\paragraph{関連問題} \hyperref[q:r01-reg-apt-08]{R01-REG-APT-08}

\par 難易度 3（校準初稿）、優先度未評価。

\subsection{令和元年度 問8}
\label{q:r01-reg-apt-08}

\noindent\textbf{原典：} sources/R01/aptitude.pdf、PDF物理9頁。

\textbf{公式正答：} 2（single）。

\paragraph{問題の要旨} 2014年行動計画のメンテナンスサイクルと施設別条件を読む。

\Needspace{4\baselineskip}\paragraph{確認した条件・記述（独自要約）}\mbox{}\par

\noindent\textbf{原典照合：} 22原典画像で25問の全五肢・数値・単位・否定を確認。\par

\noindent\textbf{対象：} 2014年行動計画のメンテナンスサイクルと施設別条件を読む。\par

\paragraph{問われる核心} 2014年行動計画のメンテナンスサイクルと施設別条件を読む。

\paragraph{必要知識} 誤り2は全分野の点検頻度・内容の全国一律化。2014年行動計画には施設別の取組があり、各施設の状態等に応じた実施と基準整備を読む。道路橋等の分野内の定期基準を、全インフラで内容まで同一という全称に変えない。

\paragraph{考え方・正答を選ぶ理由（査読待ち初稿）} 誤り2は全分野の点検頻度・内容の全国一律化。2014年行動計画には施設別の取組があり、各施設の状態等に応じた実施と基準整備を読む。道路橋等の分野内の定期基準を、全インフラで内容まで同一という全称に変えない。

\Needspace{4\baselineskip}\paragraph{選択肢の検討（査読待ち初稿）}\mbox{}\par

\noindent\textbf{1：代表的条件で適切} 点検診断・措置・記録のサイクル。施設状態を点検・診断し、必要な措置を行い記録して次周期に反映する。\par

\noindent\textbf{2：不適切} 頻度・内容を全国一律。全施設・全分野へ一つの頻度と内容を固定する記述は不適切。施設特性・状態・重要度等を考えて基準を整える。分野内の最低基準や統一した健全度尺度と混同しない。\par

\noindent\textbf{3：代表的条件で適切} 劣化損傷・構造諸元を継続収集。繰返し点検で最新の状態を記録し、過去に欠けた情報も補う。\par

\noindent\textbf{4：代表的条件で適切} ロボット・非破壊検査・ICTで高度化。診断を支える技術で効率と精度の改善を図る。新技術の適用範囲・限界を確認する。\par

\noindent\textbf{5：代表的条件で適切} 業務別能力・民間資格の認定仕組み。施設分野・業務別に必要能力を明確にし、人材の質と担い手を確保する。資格名の存在だけで全業務を担当可能とはしない。\par

\paragraph{再利用する考え方} 点検・診断・措置・記録のサイクルを再構成し、施設ごとの基準と全施設一律の主張を区別する。管理者の専門的責任と住民との協働を両立させる。

\paragraph{共通理論} \hyperref[theory:safety:maintenance]{第\ref*{theory:safety:maintenance}節}

\par 参考：\cite{eng52_infra2014}。

\paragraph{関連問題} \hyperref[q:h26-reg-apt-13]{H26-REG-APT-13}

\par 難易度 2（校準初稿）、優先度未評価。

\section{開発危険と部品原材料の免責}

\subsection{平成26年度 問7}
\label{q:h26-reg-apt-07}

\noindent\textbf{原典：} sources/H26/aptitude.pdf、PDF物理7頁。

\textbf{公式正答：} 4（single）。

\par\textbf{検討未確定：} 5の欠陥存在だけで請求できるという省略には損害と因果関係の条件を補う必要がある。公式4と初回4は保持する。

\paragraph{問題の要旨} PL法の5記述から開発危険の免責を否定する文を選ぶ。

\Needspace{4\baselineskip}\paragraph{確認した条件・記述（独自要約）}\mbox{}\par

\noindent\textbf{原典条件：} 問題幹、否定、ア等の記号順、5選択肢と個数又は表の全行を原典16画像で確認した。\par

\noindent\textbf{論理条件：} 最も不適切な単一候補を選ぶ。\par

\noindent\textbf{法の時点：} 出題時の制度を保持。現行取得基準日は2026-10-12。2020年施行の期間変更を別記する。\par

\paragraph{問われる核心} PL法の5記述から開発危険の免責を否定する文を選ぶ。

\paragraph{必要知識} 最も不適切な単一候補を選ぶ。

\paragraph{考え方・正答を選ぶ理由（査読待ち初稿）} 1〜3は製造加工動産、輸入者、通常の安全性という要件を述べる。4は引渡時の科学技術で認識不能でも免責できないとするため、4条1号と反する。5は過失立証の不要という要点を述べるが、欠陥の存在だけで全責任が成立するという独立した法律説明ではなく、損害と因果関係を前提とした省略表現として読む。

\Needspace{4\baselineskip}\paragraph{選択肢の検討（査読待ち初稿）}\mbox{}\par

\noindent\textbf{1：適切} 製造又は加工された動産であり、不動産は除く。2条1項の定義である。\par

\noindent\textbf{2：適切} 製造加工と輸入の事業者を含む。輸入者も2条3項1号の対象となる。\par

\noindent\textbf{3：適切} 通常の安全性の不足が欠陥である。安全と無関係な品質不良は区別する。\par

\noindent\textbf{4：不適切} 認識不能でも免責されない。4条1号は所定の証明による開発危険の免責を定める。\par

\noindent\textbf{5：適切} 過失立証は不要で欠陥を立証する。制度変化は正しいが、損害と因果関係を省略している。\par

\paragraph{再利用する考え方} 対象物、責任主体、欠陥・損害・因果関係、免責、期間制限を順に分けて判定する。対象外という結論から事故予防や他法の責任の不存在を導かない。

\paragraph{共通理論} \hyperref[theory:pl:scope]{第\ref*{theory:pl:scope}節}, \hyperref[theory:pl:producer]{第\ref*{theory:pl:producer}節}, \hyperref[theory:pl:defect]{第\ref*{theory:pl:defect}節}, \hyperref[theory:pl:development]{第\ref*{theory:pl:development}節}, \hyperref[theory:pl:proof]{第\ref*{theory:pl:proof}節}

\par 参考：\cite{pl_statute,pl_qa,pl_definitions2018}。

\paragraph{関連問題} \hyperref[q:h29-reg-apt-08]{H29-REG-APT-08}

\par 難易度 2（校準初稿）、優先度未評価。

\subsection{平成29年度 問8}
\label{q:h29-reg-apt-08}

\noindent\textbf{原典：} sources/H29/aptitude.pdf、PDF物理8頁。

\textbf{公式正答：} 5（single）。

\par\textbf{検討未確定：} 2は施行前の救済を709条で代表し、契約責任等を省略している。4の品質上の欠陥は安全性に関わるものと読む。文の省略範囲を査読事項に保持する。

\paragraph{問題の要旨} PL法の目的・立証・対象・誤使用・開発危険から不適切な記述を選ぶ。

\Needspace{4\baselineskip}\paragraph{確認した条件・記述（独自要約）}\mbox{}\par

\noindent\textbf{原典条件：} 問題幹、否定、ア等の記号順、5選択肢と個数又は表の全行を原典16画像で確認した。\par

\noindent\textbf{論理条件：} 最も不適切な候補。品質の欠陥は安全性に関わる範囲で理解する。\par

\noindent\textbf{法の時点：} 出題時の制度を保持。現行取得基準日は2026-10-12。2020年施行の期間変更を別記する。\par

\paragraph{問われる核心} PL法の目的・立証・対象・誤使用・開発危険から不適切な記述を選ぶ。

\paragraph{必要知識} 最も不適切な候補。品質の欠陥は安全性に関わる範囲で理解する。

\paragraph{考え方・正答を選ぶ理由（査読待ち初稿）} 1は法の目的、3は対象の定義である。2は過失から欠陥と損害の因果関係へ立証の中心が移る説明、4は通常予見される誤使用への安全確保である。5の認識不能でも責任を負うという一律の結論は開発危険の免責を落としており不適切である。

\Needspace{4\baselineskip}\paragraph{選択肢の検討（査読待ち初稿）}\mbox{}\par

\noindent\textbf{1：適切} 法は被害者保護等を目的とする。1条の趣旨である。\par

\noindent\textbf{2：適切} 欠陥と損害の関係を立証する。過失立証不要の要点は適切だが、以前の契約責任等を一括否定しない。\par

\noindent\textbf{3：適切} 対象は製造加工動産である。2条1項の定義である。\par

\noindent\textbf{4：適切} 予見可能な誤使用にも安全確保が必要。通常予見される使用形態を総合判断する。品質一般とは区別する。\par

\noindent\textbf{5：不適切} 認識不能でも責任を免れない。4条の免責条件を無視している。\par

\paragraph{再利用する考え方} 対象物、責任主体、欠陥・損害・因果関係、免責、期間制限を順に分けて判定する。対象外という結論から事故予防や他法の責任の不存在を導かない。

\paragraph{共通理論} \hyperref[theory:pl:framework]{第\ref*{theory:pl:framework}節}, \hyperref[theory:pl:proof]{第\ref*{theory:pl:proof}節}, \hyperref[theory:pl:scope]{第\ref*{theory:pl:scope}節}, \hyperref[theory:pl:defect]{第\ref*{theory:pl:defect}節}, \hyperref[theory:pl:development]{第\ref*{theory:pl:development}節}

\par 参考：\cite{pl_statute,pl_qa,pl_definitions2018}。

\paragraph{関連問題} \hyperref[q:h26-reg-apt-07]{H26-REG-APT-07}

\par 難易度 2（校準初稿）、優先度未評価。

\section{輸入者・表示製造業者・OEMの責任}

\subsection{平成28年度 問7}
\label{q:h28-reg-apt-07}

\noindent\textbf{原典：} sources/H28/aptitude.pdf、PDF物理7頁。

\textbf{公式正答：} 5（single）。

\paragraph{問題の要旨} OEM表示・無体物・不動産・未加工物の4記述を組合せに対応させる。

\Needspace{4\baselineskip}\paragraph{確認した条件・記述（独自要約）}\mbox{}\par

\noindent\textbf{原典条件：} 問題幹、否定、ア等の記号順、5選択肢と個数又は表の全行を原典16画像で確認した。\par

\noindent\textbf{論理条件：} 列順ア・イ・ウ・エ。OEMの条件付該当性。\par

\noindent\textbf{法の時点：} 出題時の制度を保持。現行取得基準日は2026-10-12。2020年施行の期間変更を別記する。\par

\paragraph{問われる核心} OEM表示・無体物・不動産・未加工物の4記述を組合せに対応させる。

\paragraph{必要知識} 列順ア・イ・ウ・エ。OEMの条件付該当性。

\paragraph{考え方・正答を選ぶ理由（査読待ち初稿）} アのOEM販売者は製造業者等に当たらないという一律除外は誤りである。製造業者と誤認させる表示等の要件を満たせば対象となる。イの無体エネルギー・ソフト・修理サービス、ウの不動産、エの未加工農林水産物と採掘したままの鉱物の除外は適切である。FTTTに対応する5を選ぶ。

\Needspace{4\baselineskip}\paragraph{選択肢の検討（査読待ち初稿）}\mbox{}\par

\noindent\textbf{1：組合せが不一致} 正・正・正・誤。各記述の判定はFTTTである。原典の列順で[1, 4]列目が異なる。\par

\noindent\textbf{2：組合せが不一致} 正・正・誤・正。各記述の判定はFTTTである。原典の列順で[1, 3]列目が異なる。\par

\noindent\textbf{3：組合せが不一致} 正・誤・正・正。各記述の判定はFTTTである。原典の列順で[1, 2]列目が異なる。\par

\noindent\textbf{4：組合せが不一致} 正・正・正・正。各記述の判定はFTTTである。原典の列順で[1]列目が異なる。\par

\noindent\textbf{5：適切な組合せ} 誤・正・正・正。各記述の判定はFTTTである。全ての列が一致する。\par

\paragraph{再利用する考え方} 対象物、責任主体、欠陥・損害・因果関係、免責、期間制限を順に分けて判定する。対象外という結論から事故予防や他法の責任の不存在を導かない。

\paragraph{共通理論} \hyperref[theory:pl:producer]{第\ref*{theory:pl:producer}節}, \hyperref[theory:pl:scope]{第\ref*{theory:pl:scope}節}, \hyperref[theory:pl:processing]{第\ref*{theory:pl:processing}節}

\par 参考：\cite{pl_statute,pl_qa,pl_definitions2018}。

\paragraph{関連問題} \hyperref[q:r06-reg-apt-14]{R06-REG-APT-14}

\par 難易度 2（校準初稿）、優先度未評価。

\section{欠陥・損害・因果関係の立証}

\subsection{平成23年度 問8}
\label{q:h23-reg-apt-08}

\noindent\textbf{原典：} sources/H23/aptitude.pdf、PDF物理9頁。

\textbf{公式正答：} 5（single）。

\paragraph{問題の要旨} PL法の立証・対象・自損・輸入者・欠陥の5記述を正誤表に対応させる。

\Needspace{4\baselineskip}\paragraph{確認した条件・記述（独自要約）}\mbox{}\par

\noindent\textbf{原典条件：} 問題幹、否定、ア等の記号順、5選択肢と個数又は表の全行を原典16画像で確認した。\par

\noindent\textbf{論理条件：} 原典列順ア・イ・ウ・エ・オ。単体ソフトと組込み製品、輸入と単なる販売、品質と安全を分ける。\par

\noindent\textbf{法の時点：} 出題時の制度を保持。現行取得基準日は2026-10-12。2020年施行の期間変更を別記する。\par

\paragraph{問われる核心} PL法の立証・対象・自損・輸入者・欠陥の5記述を正誤表に対応させる。

\paragraph{必要知識} 原典列順ア・イ・ウ・エ・オ。単体ソフトと組込み製品、輸入と単なる販売、品質と安全を分ける。

\paragraph{考え方・正答を選ぶ理由（査読待ち初稿）} アは過失の証明が法の変化であるという説明が誤り、イは電気と単体ソフトまで対象にする全称が誤りである。ウはテレビだけの損害が3条ただし書の対象外、エは輸入者を除外する点が誤りである。オは通常有すべき安全性を欠陥の基準とし、適切である。FFFFTに対応する5を選ぶ。

\Needspace{4\baselineskip}\paragraph{選択肢の検討（査読待ち初稿）}\mbox{}\par

\noindent\textbf{1：組合せが不一致} 正・正・誤・誤・正。各記述の判定はFFFFTである。原典の列順で[1, 2]列目が異なる。\par

\noindent\textbf{2：組合せが不一致} 誤・誤・正・誤・正。各記述の判定はFFFFTである。原典の列順で[3]列目が異なる。\par

\noindent\textbf{3：組合せが不一致} 正・誤・誤・誤・正。各記述の判定はFFFFTである。原典の列順で[1]列目が異なる。\par

\noindent\textbf{4：組合せが不一致} 誤・正・正・誤・誤。各記述の判定はFFFFTである。原典の列順で[2, 3, 5]列目が異なる。\par

\noindent\textbf{5：適切な組合せ} 誤・誤・誤・誤・正。各記述の判定はFFFFTである。全ての列が一致する。\par

\paragraph{再利用する考え方} 対象物、責任主体、欠陥・損害・因果関係、免責、期間制限を順に分けて判定する。対象外という結論から事故予防や他法の責任の不存在を導かない。

\paragraph{共通理論} \hyperref[theory:pl:proof]{第\ref*{theory:pl:proof}節}, \hyperref[theory:pl:scope]{第\ref*{theory:pl:scope}節}, \hyperref[theory:pl:selfdamage]{第\ref*{theory:pl:selfdamage}節}, \hyperref[theory:pl:producer]{第\ref*{theory:pl:producer}節}, \hyperref[theory:pl:defect]{第\ref*{theory:pl:defect}節}

\par 参考：\cite{pl_statute,pl_qa,pl_definitions2018}。

\paragraph{関連問題} \hyperref[q:h25-reg-apt-07]{H25-REG-APT-07}

\par 難易度 2（校準初稿）、優先度未評価。

\subsection{令和2年度 問6}
\label{q:r02-reg-apt-06}

\noindent\textbf{原典：} sources/R02/aptitude.pdf、PDF物理6頁。

\textbf{公式正答：} 5（single）。

\paragraph{問題の要旨} PL法の目的・因果関係・開発危険・民法・海外制度から最も不適切な文を選ぶ。

\Needspace{4\baselineskip}\paragraph{確認した条件・記述（独自要約）}\mbox{}\par

\noindent\textbf{原典条件：} 問題幹、否定、ア等の記号順、5選択肢と個数又は表の全行を原典16画像で確認した。\par

\noindent\textbf{論理条件：} 単一の不適切文。科学技術の認識不能は製造業者側が証明する。\par

\noindent\textbf{法の時点：} 出題時の制度を保持。現行取得基準日は2026-10-12。2020年施行の期間変更を別記する。\par

\paragraph{問われる核心} PL法の目的・因果関係・開発危険・民法・海外制度から最も不適切な文を選ぶ。

\paragraph{必要知識} 単一の不適切文。科学技術の認識不能は製造業者側が証明する。

\paragraph{考え方・正答を選ぶ理由（査読待ち初稿）} 1は目的、2は欠陥と損害の因果関係の証明、3は開発危険の免責、4は民法の補充であり適切である。5は国際的に統一された規定なので日本法だけに従えばよいと断定する。対象範囲や責任の条件は国ごとに確認すべきであり、5が不適切である。

\Needspace{4\baselineskip}\paragraph{選択肢の検討（査読待ち初稿）}\mbox{}\par

\noindent\textbf{1：適切} 被害者保護等を目的とする。1条の趣旨である。\par

\noindent\textbf{2：適切} 過失の有無にかかわらず因果関係を証明する。欠陥・損害の関係の証明は必要である。\par

\noindent\textbf{3：適切} 認識不能の開発危険を証明すれば免責。4条1号の条件を述べる。\par

\noindent\textbf{4：適切} 特段の定めがない責任は民法を適用。6条の補充規定である。\par

\noindent\textbf{5：不適切} 国際共通なので日本法だけでよい。外国の対象や要件の違いを無視している。\par

\paragraph{再利用する考え方} 対象物、責任主体、欠陥・損害・因果関係、免責、期間制限を順に分けて判定する。対象外という結論から事故予防や他法の責任の不存在を導かない。

\paragraph{共通理論} \hyperref[theory:pl:proof]{第\ref*{theory:pl:proof}節}, \hyperref[theory:pl:development]{第\ref*{theory:pl:development}節}, \hyperref[theory:pl:otherlaw]{第\ref*{theory:pl:otherlaw}節}

\par 参考：\cite{pl_statute,pl_qa,pl_definitions2018}。

\par 難易度 2（校準初稿）、優先度未評価。

\section{修理・中古品・再生品の責任境界}

\subsection{令和3年度 問8}
\label{q:r03-reg-apt-08}

\noindent\textbf{原典：} sources/R03/aptitude.pdf、PDF物理8頁。

\textbf{公式正答：} 3（single）。

\paragraph{問題の要旨} 付合・組込みソフト・再生品・修理の4記述を正誤表で判定する。

\Needspace{4\baselineskip}\paragraph{確認した条件・記述（独自要約）}\mbox{}\par

\noindent\textbf{原典条件：} 問題幹、否定、ア等の記号順、5選択肢と個数又は表の全行を原典16画像で確認した。\par

\noindent\textbf{論理条件：} 列順ア・イ・ウ・エ。修理された元の製品の対象性とは区別する。\par

\noindent\textbf{法の時点：} 出題時の制度を保持。現行取得基準日は2026-10-12。2020年施行の期間変更を別記する。\par

\paragraph{問われる核心} 付合・組込みソフト・再生品・修理の4記述を正誤表で判定する。

\paragraph{必要知識} 列順ア・イ・ウ・エ。修理された元の製品の対象性とは区別する。

\paragraph{考え方・正答を選ぶ理由（査読待ち初稿）} アは取り付け後も引渡時の動産の欠陥の責任が問題となるため正しい。イは組込みソフトの不具合と製品事故の因果関係を条件にして正しい。ウは最後の再生業者が全て責任を負うと断定するが、元の製品の欠陥と損害の因果関係があれば元の製造業者にも責任が生じ得るため誤りである。エは性質回復維持という修理サービス自体を指し正しい。TTFTが3に対応する。

\Needspace{4\baselineskip}\paragraph{選択肢の検討（査読待ち初稿）}\mbox{}\par

\noindent\textbf{1：組合せが不一致} 正・誤・正・正。各記述の判定はTTFTである。原典の列順で[2, 3]列目が異なる。\par

\noindent\textbf{2：組合せが不一致} 誤・正・正・誤。各記述の判定はTTFTである。原典の列順で[1, 3, 4]列目が異なる。\par

\noindent\textbf{3：適切な組合せ} 正・正・誤・正。各記述の判定はTTFTである。全ての列が一致する。\par

\noindent\textbf{4：組合せが不一致} 正・誤・正・誤。各記述の判定はTTFTである。原典の列順で[2, 3, 4]列目が異なる。\par

\noindent\textbf{5：組合せが不一致} 誤・正・誤・正。各記述の判定はTTFTである。原典の列順で[1]列目が異なる。\par

\paragraph{再利用する考え方} 対象物、責任主体、欠陥・損害・因果関係、免責、期間制限を順に分けて判定する。対象外という結論から事故予防や他法の責任の不存在を導かない。

\paragraph{共通理論} \hyperref[theory:pl:attachment]{第\ref*{theory:pl:attachment}節}, \hyperref[theory:pl:software]{第\ref*{theory:pl:software}節}, \hyperref[theory:pl:recycle]{第\ref*{theory:pl:recycle}節}, \hyperref[theory:pl:repair]{第\ref*{theory:pl:repair}節}

\par 参考：\cite{pl_statute,pl_qa,pl_definitions2018}。

\paragraph{関連問題} \hyperref[q:r05-reg-apt-06]{R05-REG-APT-06}

\par 難易度 2（校準初稿）、優先度未評価。

\subsection{令和6年度 問14}
\label{q:r06-reg-apt-14}

\noindent\textbf{原典：} sources/R06/aptitude.pdf、PDF物理14頁。

\textbf{公式正答：} 3（single）。

\paragraph{問題の要旨} 無体物・自損・OEM・再生品・医薬品の5記述を判定する。

\Needspace{4\baselineskip}\paragraph{確認した条件・記述（独自要約）}\mbox{}\par

\noindent\textbf{原典条件：} 問題幹、否定、ア等の記号順、5選択肢と個数又は表の全行を原典16画像で確認した。\par

\noindent\textbf{論理条件：} 列順ア・イ・ウ・エ・オ。医薬品の安全性や免責は個別に確認する。\par

\noindent\textbf{法の時点：} 出題時の制度を保持。現行取得基準日は2026-10-12。2020年施行の期間変更を別記する。\par

\paragraph{問われる核心} 無体物・自損・OEM・再生品・医薬品の5記述を判定する。

\paragraph{必要知識} 列順ア・イ・ウ・エ・オ。医薬品の安全性や免責は個別に確認する。

\paragraph{考え方・正答を選ぶ理由（査読待ち初稿）} アの電気音響サービスの対象外、イの二輪車だけの損害の除外は正しい。ウのOEM表示者は見なされることがないという一律否定は誤りである。エの製造加工された再生品の対象性は正しい。オの医薬品を一括して除外する規定はなく誤りである。TTFTFが候補3に対応する。

\Needspace{4\baselineskip}\paragraph{選択肢の検討（査読待ち初稿）}\mbox{}\par

\noindent\textbf{1：組合せが不一致} 正・誤・誤・正・誤。各記述の判定はTTFTFである。原典の列順で[2]列目が異なる。\par

\noindent\textbf{2：組合せが不一致} 誤・誤・正・誤・正。各記述の判定はTTFTFである。原典の列順で[1, 2, 3, 4, 5]列目が異なる。\par

\noindent\textbf{3：適切な組合せ} 正・正・誤・正・誤。各記述の判定はTTFTFである。全ての列が一致する。\par

\noindent\textbf{4：組合せが不一致} 正・誤・正・誤・正。各記述の判定はTTFTFである。原典の列順で[2, 3, 4, 5]列目が異なる。\par

\noindent\textbf{5：組合せが不一致} 誤・正・誤・正・誤。各記述の判定はTTFTFである。原典の列順で[1]列目が異なる。\par

\paragraph{再利用する考え方} 対象物、責任主体、欠陥・損害・因果関係、免責、期間制限を順に分けて判定する。対象外という結論から事故予防や他法の責任の不存在を導かない。

\paragraph{共通理論} \hyperref[theory:pl:scope]{第\ref*{theory:pl:scope}節}, \hyperref[theory:pl:selfdamage]{第\ref*{theory:pl:selfdamage}節}, \hyperref[theory:pl:producer]{第\ref*{theory:pl:producer}節}, \hyperref[theory:pl:recycle]{第\ref*{theory:pl:recycle}節}

\par 参考：\cite{pl_statute,pl_qa,pl_definitions2018}。

\paragraph{関連問題} \hyperref[q:h28-reg-apt-07]{H28-REG-APT-07}

\par 難易度 2（校準初稿）、優先度未評価。

\section{製造物の有体性・動産性・加工の境界}

\subsection{平成24年度 問7}
\label{q:h24-reg-apt-07}

\noindent\textbf{原典：} sources/H24/aptitude.pdf、PDF物理6頁。

\textbf{公式正答：} 4（single）。

\par\textbf{独立解答と公式正答に相違あり：} 独立解答 3。初回3・公式4。使用者が修理した製品という対象名だけでは製造物性の消滅を導けない。修理行為の非該当性又は修理起因の不具合との読み分けをAstra査読事項として保持する。

\paragraph{問題の要旨} 8対象のうちPL法の製造物に該当しない数を求める。

\Needspace{4\baselineskip}\paragraph{確認した条件・記述（独自要約）}\mbox{}\par

\noindent\textbf{原典条件：} 問題幹、否定、ア等の記号順、5選択肢と個数又は表の全行を原典16画像で確認した。\par

\noindent\textbf{論理条件：} 5候補は対象外1、2、3、4、5件。名称と責任成立の条件を分ける。\par

\noindent\textbf{法の時点：} 出題時の制度を保持。現行取得基準日は2026-10-12。2020年施行の期間変更を別記する。\par

\paragraph{問われる核心} 8対象のうちPL法の製造物に該当しない数を求める。

\paragraph{必要知識} 5候補は対象外1、2、3、4、5件。名称と責任成立の条件を分ける。

\paragraph{考え方・正答を選ぶ理由（査読待ち初稿）} アのソフト単体、ウの宅地、キの未加工トマトは対象外である。イの缶詰、オのロボット、カの有体のガスやガソリン、クの取り付け前に引き渡されたサッシ等は対象になり得る。エを修理された製品そのものと読むと製造物性が残り、対象外3件・初回3となる。公式4は対象外4件であり、エを修理行為又は修理起因の欠陥について読む場合の追加除外と整合する。消費者庁Q\&Aは中古品の製造物性を肯定し、引渡時の欠陥と後の修理を分けている。

\Needspace{4\baselineskip}\paragraph{選択肢の検討（査読待ち初稿）}\mbox{}\par

\noindent\textbf{1：個数を比較} 1件。各記述の判断理由は解法のとおりである。提示する1件は合計3件と異なる。公式は4であり、修理された製品と修理行為の読み分けを疑義に残す。\par

\noindent\textbf{2：個数を比較} 2件。各記述の判断理由は解法のとおりである。提示する2件は合計3件と異なる。公式は4であり、修理された製品と修理行為の読み分けを疑義に残す。\par

\noindent\textbf{3：個数を比較} 3件。各記述の判断理由は解法のとおりである。合計3件と一致する。公式は4であり、修理された製品と修理行為の読み分けを疑義に残す。\par

\noindent\textbf{4：個数を比較} 4件。各記述の判断理由は解法のとおりである。提示する4件は合計3件と異なる。公式は4であり、修理された製品と修理行為の読み分けを疑義に残す。\par

\noindent\textbf{5：個数を比較} 5件。各記述の判断理由は解法のとおりである。提示する5件は合計3件と異なる。公式は4であり、修理された製品と修理行為の読み分けを疑義に残す。\par

\paragraph{再利用する考え方} 対象物、責任主体、欠陥・損害・因果関係、免責、期間制限を順に分けて判定する。対象外という結論から事故予防や他法の責任の不存在を導かない。

\paragraph{共通理論} \hyperref[theory:pl:scope]{第\ref*{theory:pl:scope}節}, \hyperref[theory:pl:software]{第\ref*{theory:pl:software}節}, \hyperref[theory:pl:processing]{第\ref*{theory:pl:processing}節}, \hyperref[theory:pl:attachment]{第\ref*{theory:pl:attachment}節}, \hyperref[theory:pl:repair]{第\ref*{theory:pl:repair}節}

\par 参考：\cite{pl_statute,pl_qa,pl_definitions2018}。

\paragraph{関連問題} \hyperref[q:r01-reg-apt-03]{R01-REG-APT-03}

\par 難易度 3（校準初稿）、優先度未評価。

\subsection{平成25年度 問7}
\label{q:h25-reg-apt-07}

\noindent\textbf{原典：} sources/H25/aptitude.pdf、PDF物理8頁。

\textbf{公式正答：} 4（single）。

\par\textbf{独立解答と公式正答に相違あり：} 独立解答 5。初回5を保持する。逐条解説により冷凍冷蔵の初回分類を訂正し、5件・候補4となる。イの品質上の欠陥は安全性に関わる範囲として読む必要があり、広義の品質不良全てにPL責任を認めない。

\paragraph{問題の要旨} PL法の8記述のうち不適切な数を数える。

\Needspace{4\baselineskip}\paragraph{確認した条件・記述（独自要約）}\mbox{}\par

\noindent\textbf{原典条件：} 問題幹、否定、ア等の記号順、5選択肢と個数又は表の全行を原典16画像で確認した。\par

\noindent\textbf{論理条件：} 候補は不適切2、3、4、5、6件。冷凍冷蔵とハム等の加工を対比する。\par

\noindent\textbf{法の時点：} 出題時の制度を保持。現行取得基準日は2026-10-12。2020年施行の期間変更を別記する。\par

\paragraph{問われる核心} PL法の8記述のうち不適切な数を数える。

\paragraph{必要知識} 候補は不適切2、3、4、5、6件。冷凍冷蔵とハム等の加工を対比する。

\paragraph{考え方・正答を選ぶ理由（査読待ち初稿）} アは電気やソフトまで対象にする点、ウは開発危険免責を否定する点、エは過失立証の制度説明、カは製品自体だけの損害を含める点、キは輸入者除外が不適切である。イは予見可能な誤使用への安全確保、クは安全と無関係の品質不具合の除外を述べる。オの単なる冷凍冷蔵を未加工とする説明は、逐条解説の物理3頁を画像で確認して適切と訂正した。初回はオも誤りとして6件・5であったが、訂正後は5件・4となる。公式4との一致とは独立した一次資料上の訂正である。

\Needspace{4\baselineskip}\paragraph{選択肢の検討（査読待ち初稿）}\mbox{}\par

\noindent\textbf{1：個数を比較} 2件。各記述の判断理由は解法のとおりである。提示する2件は合計5件と異なる。初回分類訂正後の合計であり初回5を保持する。\par

\noindent\textbf{2：個数を比較} 3件。各記述の判断理由は解法のとおりである。提示する3件は合計5件と異なる。初回分類訂正後の合計であり初回5を保持する。\par

\noindent\textbf{3：個数を比較} 4件。各記述の判断理由は解法のとおりである。提示する4件は合計5件と異なる。初回分類訂正後の合計であり初回5を保持する。\par

\noindent\textbf{4：個数を比較} 5件。各記述の判断理由は解法のとおりである。合計5件と一致する。初回分類訂正後の合計であり初回5を保持する。\par

\noindent\textbf{5：個数を比較} 6件。各記述の判断理由は解法のとおりである。提示する6件は合計5件と異なる。初回分類訂正後の合計であり初回5を保持する。\par

\paragraph{再利用する考え方} 対象物、責任主体、欠陥・損害・因果関係、免責、期間制限を順に分けて判定する。対象外という結論から事故予防や他法の責任の不存在を導かない。

\paragraph{共通理論} \hyperref[theory:pl:proof]{第\ref*{theory:pl:proof}節}, \hyperref[theory:pl:scope]{第\ref*{theory:pl:scope}節}, \hyperref[theory:pl:processing]{第\ref*{theory:pl:processing}節}, \hyperref[theory:pl:development]{第\ref*{theory:pl:development}節}, \hyperref[theory:pl:selfdamage]{第\ref*{theory:pl:selfdamage}節}, \hyperref[theory:pl:producer]{第\ref*{theory:pl:producer}節}

\par 参考：\cite{pl_statute,pl_qa,pl_definitions2018}。

\paragraph{関連問題} \hyperref[q:h23-reg-apt-08]{H23-REG-APT-08}

\par 難易度 3（校準初稿）、優先度未評価。

\subsection{令和元年度 問3}
\label{q:r01-reg-apt-03}

\noindent\textbf{原典：} sources/R01/aptitude.pdf、PDF物理4頁。

\textbf{公式正答：} 4（single）。

\paragraph{問題の要旨} 8事故でPL法の賠償責任に該当しない事例の数を求める。

\Needspace{4\baselineskip}\paragraph{確認した条件・記述（独自要約）}\mbox{}\par

\noindent\textbf{原典条件：} 問題幹、否定、ア等の記号順、5選択肢と個数又は表の全行を原典16画像で確認した。\par

\noindent\textbf{論理条件：} 対象外数候補8、7、6、5、4。製造欠陥と後の保守ミスを分ける。\par

\noindent\textbf{法の時点：} 出題時の制度を保持。現行取得基準日は2026-10-12。2020年施行の期間変更を別記する。\par

\paragraph{問われる核心} 8事故でPL法の賠償責任に該当しない事例の数を求める。

\paragraph{必要知識} 対象外数候補8、7、6、5、4。製造欠陥と後の保守ミスを分ける。

\paragraph{考え方・正答を選ぶ理由（査読待ち初稿）} アの輸入車、ウの電動シャッター、カの組込みソフトによるロボットは、製造時の欠陥と負傷の因果関係を前提にPL責任の対象となる。イの宅地、エの未加工鶏卵、オの保守サービス、キの電力、クの未加工養殖魚は対象外で5件である。候補4を選ぶ。保守ミスの責任や食品の安全確保が消えるという意味ではない。

\Needspace{4\baselineskip}\paragraph{選択肢の検討（査読待ち初稿）}\mbox{}\par

\noindent\textbf{1：個数を比較} 8件。各記述の判断理由は解法のとおりである。提示する8件は合計5件と異なる。\par

\noindent\textbf{2：個数を比較} 7件。各記述の判断理由は解法のとおりである。提示する7件は合計5件と異なる。\par

\noindent\textbf{3：個数を比較} 6件。各記述の判断理由は解法のとおりである。提示する6件は合計5件と異なる。\par

\noindent\textbf{4：個数を比較} 5件。各記述の判断理由は解法のとおりである。合計5件と一致する。\par

\noindent\textbf{5：個数を比較} 4件。各記述の判断理由は解法のとおりである。提示する4件は合計5件と異なる。\par

\paragraph{再利用する考え方} 対象物、責任主体、欠陥・損害・因果関係、免責、期間制限を順に分けて判定する。対象外という結論から事故予防や他法の責任の不存在を導かない。

\paragraph{共通理論} \hyperref[theory:pl:scope]{第\ref*{theory:pl:scope}節}, \hyperref[theory:pl:software]{第\ref*{theory:pl:software}節}, \hyperref[theory:pl:attachment]{第\ref*{theory:pl:attachment}節}, \hyperref[theory:pl:processing]{第\ref*{theory:pl:processing}節}, \hyperref[theory:pl:repair]{第\ref*{theory:pl:repair}節}

\par 参考：\cite{pl_statute,pl_qa,pl_definitions2018}。

\paragraph{関連問題} \hyperref[q:h24-reg-apt-07]{H24-REG-APT-07}, \hyperref[q:r04-reg-apt-11]{R04-REG-APT-11}

\par 難易度 3（校準初稿）、優先度未評価。

\subsection{令和4年度 問11}
\label{q:r04-reg-apt-11}

\noindent\textbf{原典：} sources/R04/aptitude.pdf、PDF物理10頁。

\textbf{公式正答：} 4（single）。

\paragraph{問題の要旨} 時効期限内の8事故からPL法の賠償責任に該当しない数を求める。

\Needspace{4\baselineskip}\paragraph{確認した条件・記述（独自要約）}\mbox{}\par

\noindent\textbf{原典条件：} 問題幹、否定、ア等の記号順、5選択肢と個数又は表の全行を原典16画像で確認した。\par

\noindent\textbf{論理条件：} 全事例は時効期限内。対象外数候補1、2、3、4、5。\par

\noindent\textbf{法の時点：} 出題時の制度を保持。現行取得基準日は2026-10-12。2020年施行の期間変更を別記する。\par

\paragraph{問われる核心} 時効期限内の8事故からPL法の賠償責任に該当しない数を求める。

\paragraph{必要知識} 全事例は時効期限内。対象外数候補1、2、3、4、5。

\paragraph{考え方・正答を選ぶ理由（査読待ち初稿）} アの冷蔵庫から住宅への延焼、ウの製造欠陥のあるエスカレータ、オの組込みソフトのロボット、キの輸入加工品の生ハムは対象となる。イの土地建売住宅の施工不良、エの電力、カの未加工自然養殖活魚、クの後の保守作業ミスは対象外で4件となり候補4である。時効期限内という明示条件により期間判定をここで追加しない。

\Needspace{4\baselineskip}\paragraph{選択肢の検討（査読待ち初稿）}\mbox{}\par

\noindent\textbf{1：個数を比較} 1件。各記述の判断理由は解法のとおりである。提示する1件は合計4件と異なる。\par

\noindent\textbf{2：個数を比較} 2件。各記述の判断理由は解法のとおりである。提示する2件は合計4件と異なる。\par

\noindent\textbf{3：個数を比較} 3件。各記述の判断理由は解法のとおりである。提示する3件は合計4件と異なる。\par

\noindent\textbf{4：個数を比較} 4件。各記述の判断理由は解法のとおりである。合計4件と一致する。\par

\noindent\textbf{5：個数を比較} 5件。各記述の判断理由は解法のとおりである。提示する5件は合計4件と異なる。\par

\paragraph{再利用する考え方} 対象物、責任主体、欠陥・損害・因果関係、免責、期間制限を順に分けて判定する。対象外という結論から事故予防や他法の責任の不存在を導かない。

\paragraph{共通理論} \hyperref[theory:pl:scope]{第\ref*{theory:pl:scope}節}, \hyperref[theory:pl:attachment]{第\ref*{theory:pl:attachment}節}, \hyperref[theory:pl:software]{第\ref*{theory:pl:software}節}, \hyperref[theory:pl:processing]{第\ref*{theory:pl:processing}節}, \hyperref[theory:pl:producer]{第\ref*{theory:pl:producer}節}, \hyperref[theory:pl:repair]{第\ref*{theory:pl:repair}節}

\par 参考：\cite{pl_statute,pl_qa,pl_definitions2018}。

\paragraph{関連問題} \hyperref[q:r01-reg-apt-03]{R01-REG-APT-03}

\par 難易度 3（校準初稿）、優先度未評価。

\subsection{令和5年度 問6}
\label{q:r05-reg-apt-06}

\noindent\textbf{原典：} sources/R05/aptitude.pdf、PDF物理6頁。

\textbf{公式正答：} 3（single）。

\paragraph{問題の要旨} 付合・組込みソフト・原子力・修理・海外共通化の5記述を判定する。

\Needspace{4\baselineskip}\paragraph{確認した条件・記述（独自要約）}\mbox{}\par

\noindent\textbf{原典条件：} 問題幹、否定、ア等の記号順、5選択肢と個数又は表の全行を原典16画像で確認した。\par

\noindent\textbf{論理条件：} 列順ア・イ・ウ・エ・オ。特別法と修理サービスを区別。\par

\noindent\textbf{法の時点：} 出題時の制度を保持。現行取得基準日は2026-10-12。2020年施行の期間変更を別記する。\par

\paragraph{問われる核心} 付合・組込みソフト・原子力・修理・海外共通化の5記述を判定する。

\paragraph{必要知識} 列順ア・イ・ウ・エ・オ。特別法と修理サービスを区別。

\paragraph{考え方・正答を選ぶ理由（査読待ち初稿）} アのエスカレータはいかなる場合も対象外という断定は誤りである。イの組込み製品の因果関係、ウの原子力特別法、エの性質回復維持という修理サービスの除外は正しい。オの日本法だけで海外対応できるという説明は誤りである。FTTTFが候補3に対応する。エは修理された元の製品自体の対象性を一括して否定する文として読まない。

\Needspace{4\baselineskip}\paragraph{選択肢の検討（査読待ち初稿）}\mbox{}\par

\noindent\textbf{1：組合せが不一致} 正・誤・正・正・誤。各記述の判定はFTTTFである。原典の列順で[1, 2]列目が異なる。\par

\noindent\textbf{2：組合せが不一致} 正・正・誤・誤・正。各記述の判定はFTTTFである。原典の列順で[1, 3, 4, 5]列目が異なる。\par

\noindent\textbf{3：適切な組合せ} 誤・正・正・正・誤。各記述の判定はFTTTFである。全ての列が一致する。\par

\noindent\textbf{4：組合せが不一致} 誤・誤・正・正・誤。各記述の判定はFTTTFである。原典の列順で[2]列目が異なる。\par

\noindent\textbf{5：組合せが不一致} 誤・誤・誤・誤・正。各記述の判定はFTTTFである。原典の列順で[2, 3, 4, 5]列目が異なる。\par

\paragraph{再利用する考え方} 対象物、責任主体、欠陥・損害・因果関係、免責、期間制限を順に分けて判定する。対象外という結論から事故予防や他法の責任の不存在を導かない。

\paragraph{共通理論} \hyperref[theory:pl:attachment]{第\ref*{theory:pl:attachment}節}, \hyperref[theory:pl:software]{第\ref*{theory:pl:software}節}, \hyperref[theory:pl:otherlaw]{第\ref*{theory:pl:otherlaw}節}, \hyperref[theory:pl:repair]{第\ref*{theory:pl:repair}節}

\par 参考：\cite{pl_statute,pl_qa,pl_definitions2018,pl_nuclear_statute}。

\paragraph{関連問題} \hyperref[q:r01-ret-apt-07]{R01-RET-APT-07}, \hyperref[q:r03-reg-apt-08]{R03-REG-APT-08}

\par 難易度 2（校準初稿）、優先度未評価。

\section{製造物責任の期間制限と法改正}

\subsection{平成30年度 問9}
\label{q:h30-reg-apt-09}

\noindent\textbf{原典：} sources/H30/aptitude.pdf、PDF物理9頁。

\textbf{公式正答：} 3（single）。

\par\textbf{検討未確定：} エは購入日と製造業者等の引渡し日の関係を明示しない。一般のテレビ火災と欠陥だけで請求できるとの断定として不適切と判定し、期間特則や中断等の個別事情の不存在まで推測しない。

\paragraph{問題の要旨} PL法の6記述で、予見可能な誤使用と安全性の定義以外の誤りを数える。

\Needspace{4\baselineskip}\paragraph{確認した条件・記述（独自要約）}\mbox{}\par

\noindent\textbf{原典条件：} 問題幹、否定、ア等の記号順、5選択肢と個数又は表の全行を原典16画像で確認した。\par

\noindent\textbf{論理条件：} 5候補は不適切2、3、4、5、6件。出題は2020年改正前。\par

\noindent\textbf{法の時点：} 出題時の制度を保持。現行取得基準日は2026-10-12。2020年施行の期間変更を別記する。\par

\paragraph{問われる核心} PL法の6記述で、予見可能な誤使用と安全性の定義以外の誤りを数える。

\paragraph{必要知識} 5候補は不適切2、3、4、5、6件。出題は2020年改正前。

\paragraph{考え方・正答を選ぶ理由（査読待ち初稿）} アとカは適切である。イは開発危険、ウは過失立証、オは輸入者について誤る。エは通常のテレビ火災に購入後10年を過ぎても欠陥の証明だけで請求できると断定し、引渡し起算の期間制限と他の成立要件を無視する。したがって不適切はイ・ウ・エ・オの4件で候補3となる。購入と引渡しは同一日と決め付けず、潜伏性損害の例外とも分ける。

\Needspace{4\baselineskip}\paragraph{選択肢の検討（査読待ち初稿）}\mbox{}\par

\noindent\textbf{1：個数を比較} 2件。各記述の判断理由は解法のとおりである。提示する2件は合計4件と異なる。\par

\noindent\textbf{2：個数を比較} 3件。各記述の判断理由は解法のとおりである。提示する3件は合計4件と異なる。\par

\noindent\textbf{3：個数を比較} 4件。各記述の判断理由は解法のとおりである。合計4件と一致する。\par

\noindent\textbf{4：個数を比較} 5件。各記述の判断理由は解法のとおりである。提示する5件は合計4件と異なる。\par

\noindent\textbf{5：個数を比較} 6件。各記述の判断理由は解法のとおりである。提示する6件は合計4件と異なる。\par

\paragraph{再利用する考え方} 対象物、責任主体、欠陥・損害・因果関係、免責、期間制限を順に分けて判定する。対象外という結論から事故予防や他法の責任の不存在を導かない。

\paragraph{共通理論} \hyperref[theory:pl:defect]{第\ref*{theory:pl:defect}節}, \hyperref[theory:pl:development]{第\ref*{theory:pl:development}節}, \hyperref[theory:pl:proof]{第\ref*{theory:pl:proof}節}, \hyperref[theory:pl:time]{第\ref*{theory:pl:time}節}, \hyperref[theory:pl:producer]{第\ref*{theory:pl:producer}節}

\par 参考：\cite{pl_statute,pl_qa,pl_transition2020}。

\paragraph{関連問題} \hyperref[q:r07-reg-apt-12]{R07-REG-APT-12}

\par 難易度 3（校準初稿）、優先度未評価。

\subsection{令和元年度再試験 問7}
\label{q:r01-ret-apt-07}

\noindent\textbf{原典：} sources/R01\_retest/aptitude.pdf、PDF物理7頁。

\textbf{公式正答：} 3（single）。

\paragraph{問題の要旨} PL法の不動産例外・欠陥・時効・原子力・民法の5記述から正しい数を求める。

\Needspace{4\baselineskip}\paragraph{確認した条件・記述（独自要約）}\mbox{}\par

\noindent\textbf{原典条件：} 問題幹、否定、ア等の記号順、5選択肢と個数又は表の全行を原典16画像で確認した。\par

\noindent\textbf{論理条件：} 正しい数候補1、2、3、4、5。通常試験と再試験を別IDで保持。\par

\noindent\textbf{法の時点：} 出題時の制度を保持。現行取得基準日は2026-10-12。2020年施行の期間変更を別記する。\par

\paragraph{問われる核心} PL法の不動産例外・欠陥・時効・原子力・民法の5記述から正しい数を求める。

\paragraph{必要知識} 正しい数候補1、2、3、4、5。通常試験と再試験を別IDで保持。

\paragraph{考え方・正答を選ぶ理由（査読待ち初稿）} アの戸建て住宅耐震違反をPL法の例外とする説明は誤りであり、ウの時効なしも誤りである。イの欠陥の定義、エの原子力損害へのPL法不適用、オの民法の補充適用は適切である。正しい3件が候補3に対応する。原子力損害は特別法4条3項を確認し、住宅の他法による責任とは分ける。

\Needspace{4\baselineskip}\paragraph{選択肢の検討（査読待ち初稿）}\mbox{}\par

\noindent\textbf{1：個数を比較} 1件。各記述の判断理由は解法のとおりである。提示する1件は合計3件と異なる。\par

\noindent\textbf{2：個数を比較} 2件。各記述の判断理由は解法のとおりである。提示する2件は合計3件と異なる。\par

\noindent\textbf{3：個数を比較} 3件。各記述の判断理由は解法のとおりである。合計3件と一致する。\par

\noindent\textbf{4：個数を比較} 4件。各記述の判断理由は解法のとおりである。提示する4件は合計3件と異なる。\par

\noindent\textbf{5：個数を比較} 5件。各記述の判断理由は解法のとおりである。提示する5件は合計3件と異なる。\par

\paragraph{再利用する考え方} 対象物、責任主体、欠陥・損害・因果関係、免責、期間制限を順に分けて判定する。対象外という結論から事故予防や他法の責任の不存在を導かない。

\paragraph{共通理論} \hyperref[theory:pl:scope]{第\ref*{theory:pl:scope}節}, \hyperref[theory:pl:defect]{第\ref*{theory:pl:defect}節}, \hyperref[theory:pl:time]{第\ref*{theory:pl:time}節}, \hyperref[theory:pl:otherlaw]{第\ref*{theory:pl:otherlaw}節}

\par 参考：\cite{pl_statute,pl_qa,pl_nuclear_statute,pl_transition2020}。

\paragraph{関連問題} \hyperref[q:r05-reg-apt-06]{R05-REG-APT-06}

\par 難易度 3（校準初稿）、優先度未評価。

\subsection{令和7年度 問12}
\label{q:r07-reg-apt-12}

\noindent\textbf{原典：} sources/R07/aptitude.pdf、PDF物理12頁。

\textbf{公式正答：} 4（single）。

\paragraph{問題の要旨} PL法の対象・欠陥・OEM・時効・民法の5文から不適切なものを選ぶ。

\Needspace{4\baselineskip}\paragraph{確認した条件・記述（独自要約）}\mbox{}\par

\noindent\textbf{原典条件：} 問題幹、否定、ア等の記号順、5選択肢と個数又は表の全行を原典16画像で確認した。\par

\noindent\textbf{論理条件：} 最も不適切な単一文。OEMは製造業者表示とみなされる場合という条件付き。\par

\noindent\textbf{法の時点：} 出題時の制度を保持。現行取得基準日は2026-10-12。2020年施行の期間変更を別記する。\par

\paragraph{問われる核心} PL法の対象・欠陥・OEM・時効・民法の5文から不適切なものを選ぶ。

\paragraph{必要知識} 最も不適切な単一文。OEMは製造業者表示とみなされる場合という条件付き。

\paragraph{考え方・正答を選ぶ理由（査読待ち初稿）} 1は無体物とサービスの除外、2は安全性と単なる品質の区別、3は表示製造業者の条件付き該当性、5は民法補充であり適切である。4は消費者保護を理由に時効がないと断定するが、5条が期間制限を定めるため不適切である。初回4と公式4が一致する。

\Needspace{4\baselineskip}\paragraph{選択肢の検討（査読待ち初稿）}\mbox{}\par

\noindent\textbf{1：適切} 製造加工動産が対象で無体物等は除く。2条1項の対象条件に対応する。\par

\noindent\textbf{2：適切} 通常の安全性の欠如を欠陥とする。品質一般から安全性に関係するものを区別する。\par

\noindent\textbf{3：適切} 製造業者表示とみなされるOEM先は対象。2条3項の表示者に該当する条件を明示している。\par

\noindent\textbf{4：不適切} 請求権には時効がない。5条に短期と長期の期間制限がある。\par

\noindent\textbf{5：適切} 特段の規定がない責任は民法による。6条が補充適用を定める。\par

\paragraph{再利用する考え方} 対象物、責任主体、欠陥・損害・因果関係、免責、期間制限を順に分けて判定する。対象外という結論から事故予防や他法の責任の不存在を導かない。

\paragraph{共通理論} \hyperref[theory:pl:scope]{第\ref*{theory:pl:scope}節}, \hyperref[theory:pl:defect]{第\ref*{theory:pl:defect}節}, \hyperref[theory:pl:producer]{第\ref*{theory:pl:producer}節}, \hyperref[theory:pl:time]{第\ref*{theory:pl:time}節}, \hyperref[theory:pl:otherlaw]{第\ref*{theory:pl:otherlaw}節}

\par 参考：\cite{pl_statute,pl_qa,pl_transition2020}。

\paragraph{関連問題} \hyperref[q:h30-reg-apt-09]{H30-REG-APT-09}

\par 難易度 2（校準初稿）、優先度未評価。

\section{製品安全の社会的責任と予防}

\subsection{令和2年度 問12}
\label{q:r02-reg-apt-12}

\noindent\textbf{原典：} sources/R02/aptitude.pdf、PDF物理12頁。

\textbf{公式正答：} 1（single）。

\par\textbf{検討未確定：} ISO26000原規格全文と経済産業省研修PDF全文は未取得。問題提示定義と全7取組は原典確認し、索引による補助確認の範囲を超えて規格版照合済みとは扱わない。

\paragraph{問題の要旨} 製品安全の社会的責任に関する7取組から不適切な数を求める。

\Needspace{4\baselineskip}\paragraph{確認した条件・記述（独自要約）}\mbox{}\par

\noindent\textbf{原典条件：} 問題幹、否定、ア等の記号順、5選択肢と個数又は表の全行を原典16画像で確認した。\par

\noindent\textbf{論理条件：} 問題提示のISO26000の定義を条件として用いる。候補0、1、2、3、4件。\par

\noindent\textbf{法の時点：} 出題時の制度を保持。現行取得基準日は2026-10-12。2020年施行の期間変更を別記する。\par

\paragraph{問われる核心} 製品安全の社会的責任に関する7取組から不適切な数を求める。

\paragraph{必要知識} 問題提示のISO26000の定義を条件として用いる。候補0、1、2、3、4件。

\paragraph{考え方・正答を選ぶ理由（査読待ち初稿）} アの法令順守を超えるリスク低減、イの消費者期待を踏まえた基準、ウの事故予防、エの対話、オの将来と弱者への配慮、カの緊急時対応、キの苦情解決手段は、いずれも問題が提示する社会的責任の定義と整合する。賠償責任を負わないことだけへ活動を限定しない。不適切は0件で候補1となる。

\Needspace{4\baselineskip}\paragraph{選択肢の検討（査読待ち初稿）}\mbox{}\par

\noindent\textbf{1：個数を比較} 0件。各記述の判断理由は解法のとおりである。合計0件と一致する。\par

\noindent\textbf{2：個数を比較} 1件。各記述の判断理由は解法のとおりである。提示する1件は合計0件と異なる。\par

\noindent\textbf{3：個数を比較} 2件。各記述の判断理由は解法のとおりである。提示する2件は合計0件と異なる。\par

\noindent\textbf{4：個数を比較} 3件。各記述の判断理由は解法のとおりである。提示する3件は合計0件と異なる。\par

\noindent\textbf{5：個数を比較} 4件。各記述の判断理由は解法のとおりである。提示する4件は合計0件と異なる。\par

\paragraph{再利用する考え方} 対象物、責任主体、欠陥・損害・因果関係、免責、期間制限を順に分けて判定する。対象外という結論から事故予防や他法の責任の不存在を導かない。

\paragraph{共通理論} \hyperref[theory:pl:csr]{第\ref*{theory:pl:csr}節}, \hyperref[theory:pl:framework]{第\ref*{theory:pl:framework}節}

\par 参考：\cite{pl_product_csr2016}。

\par 難易度 3（校準初稿）、優先度未評価。

\section{製品の警告・使用情報と予見可能な誤使用}

\subsection{平成25年度 問9}
\label{q:h25-reg-apt-09}

\noindent\textbf{原典：} sources/H25/aptitude.pdf、PDF物理10頁。

\textbf{公式正答：} 5（single）。

\par\textbf{検討未確定：} 確認できた家電協会公開版は第5版2020で、出題2013時点の版の本文・国外適用条項は未照合。Guide51公開全文も取得不能。公式5と独立5を保持。

\paragraph{問題の要旨} 家電の警告表示の対象段階、法令優先、媒体、適用地域と予見可能な誤使用を整理する。

\Needspace{4\baselineskip}\paragraph{確認した条件・記述（独自要約）}\mbox{}\par

\noindent\textbf{判断方向：} 各肢の理由に記録。組合せ・個数問題では記述の真偽と回答番号を区別する。\par

\noindent\textbf{原典照合：} 全5候補、表の列順、否定、条件、規格番号、法令時点を23原典画像で照合。\par

\paragraph{問われる核心} 家電の警告表示の対象段階、法令優先、媒体、適用地域と予見可能な誤使用を整理する。

\paragraph{必要知識} ライフサイクル、警告の段階、法令、媒体、適用範囲、予見可能な誤使用。

\paragraph{考え方・正答を選ぶ理由（査読待ち初稿）} 最も不適切な5を選ぶ。誤使用という分類だけで危害の説明を排除しない。警告は設計・防護を補完するものであり、予見可能な誤使用と全く予測できない使用を区別する。4の指針の適用範囲と現地の安全責務も分ける。

\Needspace{4\baselineskip}\paragraph{選択肢の検討（査読待ち初稿）}\mbox{}\par

\noindent\textbf{1：対象段階として適切} 購入から廃棄、保守保管を対象。。使用だけでなく運搬、保守、保管、廃棄で生じる危険も検討する。\par

\noindent\textbf{2：法令優先として適切} 法定表示は関係法令に従う。。自主的ガイドラインだけで法定の表示義務を上書きできない。\par

\noindent\textbf{3：媒体として適切} 本体以外の説明書・カタログ・包装にも表示。。危険に接する場面で届く媒体を組み合わせる。本体に必要な表示を他媒体だけで代替できるという意味ではない。\par

\noindent\textbf{4：版・適用条項要確認} 国内仕様の指針なので国外では本指針の表示不要。。国内用指針の適用範囲と輸出先の法令・規格に従う義務は別である。国外なら一切の警告が不要とはしない。\par

\noindent\textbf{5：不適切} 誤使用で生じる危害は警告不要。。合理的に予見できる誤使用にも対応する。製造者の意図した使用以外だから無条件で表示不要とはいえない。\par

\paragraph{再利用する考え方} 危険源、曝露、損害、評価基準、対策、残留リスクを分ける。全称・必須・予算だけ・優先順位なしという限定を検討し、出題当時と現在の規定を混同しない。

\paragraph{共通理論} \hyperref[theory:safety:warnings]{第\ref*{theory:safety:warnings}節}, \hyperref[theory:safety:hierarchy]{第\ref*{theory:safety:hierarchy}節}

\par 参考：\cite{saf_aeha,saf_guide51}。

\par 難易度 2（校準初稿）、優先度未評価。

\subsection{令和3年度 問10}
\label{q:r03-reg-apt-10}

\noindent\textbf{原典：} sources/R03/aptitude.pdf、PDF物理10頁。

\textbf{公式正答：} 3（single）。

\par\textbf{検討未確定：} ウの輸入国が未指定で英語が公用語又は適切な場合がある。Guide51の言語条項全文を未取得。ウの無条件な誤り認定には疑義を残し公式3・初回3を保持。

\paragraph{問題の要旨} 設計軽量化の優先、ALARPの費用条件、使用国で理解できる警告と保守情報を検討する。

\Needspace{4\baselineskip}\paragraph{確認した条件・記述（独自要約）}\mbox{}\par

\noindent\textbf{判断方向：} 各肢の理由に記録。組合せ・個数問題では記述の真偽と回答番号を区別する。\par

\noindent\textbf{原典照合：} 全5候補、表の列順、否定、条件、規格番号、法令時点を23原典画像で照合。\par

\paragraph{問われる核心} 設計軽量化の優先、ALARPの費用条件、使用国で理解できる警告と保守情報を検討する。

\paragraph{必要知識} 本質設計、予算と著しい不均衡、使用先言語、警告、保守。

\paragraph{考え方・正答を選ぶ理由（査読待ち初稿）} 初回は3個の3で公式と一致する。アは軽量化という有効な上位策を意匠のため省略し、イは予算との比較を著しい費用・効果不均衡と取り違える。ウは製造国の言語と英語だけで輸入先の利用者へ必ず理解可能とは保証できない。ただし輸入国が特定されず、英語が適切な場合まで誤りと断定できるかは追加査読へ残す。

\Needspace{4\baselineskip}\paragraph{選択肢の検討（査読待ち初稿）}\mbox{}\par

\noindent\textbf{1：不足} 誤り1個。アは意匠優先で設計対策を省略、イは僅かな予算超過だけで省略し、少なくとも二つ誤る。\par

\noindent\textbf{2：不足} 誤り2個。アとイだけでは、ウの使用国への情報という問題を落とす。\par

\noindent\textbf{3：初回・公式一致、ウの条件疑義} 誤り3個。初回はア・イ・ウが誤りで3個。エは危険源と危害・結果、オは意思決定と誤使用回避、カは保守方法と訓練・保護具の情報として適切。\par

\noindent\textbf{4：過剰} 誤り4個。エ・オ・カに第四の明確な誤りはない。\par

\noindent\textbf{5：過剰} 誤り5個。エ・オ・カを一律に否定する根拠はない。\par

\paragraph{再利用する考え方} 危険源、曝露、損害、評価基準、対策、残留リスクを分ける。全称・必須・予算だけ・優先順位なしという限定を検討し、出題当時と現在の規定を混同しない。

\paragraph{共通理論} \hyperref[theory:safety:warnings]{第\ref*{theory:safety:warnings}節}, \hyperref[theory:safety:hierarchy]{第\ref*{theory:safety:hierarchy}節}, \hyperref[theory:safety:alarp]{第\ref*{theory:safety:alarp}節}

\par 参考：\cite{saf_ra2006,saf_guide51,saf_parliament2006}。

\paragraph{関連問題} \hyperref[q:r05-reg-apt-14]{R05-REG-APT-14}

\par 難易度 3（校準初稿）、優先度未評価。

\section{ALARAと防護の経済社会的最適化}

\subsection{平成26年度 問15}
\label{q:h26-reg-apt-15}

\noindent\textbf{原典：} sources/H26/aptitude.pdf、PDF物理15頁。

\textbf{公式正答：} 3（single）。

\par\textbf{検討未確定：} 2008年方針の全文は取得できず2023年一次資料の回顧と方針説明で確認。肢5の当時の詳細文言全体の照合は保留。肢1は防護モデルの仮定に限定。

\paragraph{問題の要旨} 放射線防護と食品汚染物質のALARAを、経済・社会条件を伴う最適化として理解する。

\Needspace{4\baselineskip}\paragraph{確認した条件・記述（独自要約）}\mbox{}\par

\noindent\textbf{判断方向：} 各肢の理由に記録。組合せ・個数問題では記述の真偽と回答番号を区別する。\par

\noindent\textbf{原典照合：} 全5候補、表の列順、否定、条件、規格番号、法令時点を23原典画像で照合。\par

\paragraph{問われる核心} 放射線防護と食品汚染物質のALARAを、経済・社会条件を伴う最適化として理解する。

\paragraph{必要知識} LNT防護モデル、確率的影響、最適化、経済社会要因、国内曝露と国際食品基準。

\paragraph{考え方・正答を選ぶ理由（査読待ち初稿）} 最も不適切な3を選ぶ。合理的達成可能性に含まれる経済・社会条件を一律に除外する点が2の原理と矛盾する。1の線量と確率的影響の関係は防護上の仮定として扱い、組織反応等のしきい値を混同しない。

\Needspace{4\baselineskip}\paragraph{選択肢の検討（査読待ち初稿）}\mbox{}\par

\noindent\textbf{1：適切（条件内）} ICRP2007、しきい値なしの確率的影響を前提に最適化。。低線量でも線量に応じた確率的影響を仮定する防護モデルである。あらゆる放射線影響にしきい値がないとの実験的確定ではない。\par

\noindent\textbf{2：適切（条件内）} 経済的・社会的要因を考慮。。ICRP103第203項は曝露確率・人数・線量について経済社会条件を含め最適化する。\par

\noindent\textbf{3：不適切} 自然放射線地域では費用・生活制約を考えず常に低減優先。。自然由来や既存の曝露でも最適化を行うが、経済・社会条件を除外して常に最小線量を優先する合意ではない。\par

\noindent\textbf{4：適切（条件内）} 食品汚染物質の低減・基準設定にも適用。。厚労省2023資料は2008年方針以来、食品汚染物質でALARAを考慮してきたと説明する。\par

\noindent\textbf{5：適切（条件内）} 国内の実態・摂取を踏まえコーデックス規格を検討。。国内の汚染分布・摂取と健康指標を検討するので、国際基準を無検討に採用する制度ではない。原典の当時の表現として保存する。\par

\paragraph{再利用する考え方} 危険源、曝露、損害、評価基準、対策、残留リスクを分ける。全称・必須・予算だけ・優先順位なしという限定を検討し、出題当時と現在の規定を混同しない。

\paragraph{共通理論} \hyperref[theory:safety:alara]{第\ref*{theory:safety:alara}節}, \hyperref[theory:safety:quantification]{第\ref*{theory:safety:quantification}節}

\par 参考：\cite{saf_icrp103,saf_icrpaedia,saf_food2023}。

\par 難易度 2（校準初稿）、優先度未評価。

\section{安全の許容性・残留リスクと社会条件}

\subsection{平成23年度 問5}
\label{q:h23-reg-apt-05}

\noindent\textbf{原典：} sources/H23/aptitude.pdf、PDF物理6頁。

\textbf{公式正答：} 4（single）。

\par\textbf{検討未確定：} 肢1・2の日本と国際の安全文化の全称的対比の根拠未確認。Guide51公開本文取得は403で未読、公式4への一致で全肢の一般性を確定しない。

\paragraph{問題の要旨} 安全の許容基準の社会・技術依存性と、設計から情報提供までの対策順位を判定する。

\Needspace{4\baselineskip}\paragraph{確認した条件・記述（独自要約）}\mbox{}\par

\noindent\textbf{判断方向：} 各肢の理由に記録。組合せ・個数問題では記述の真偽と回答番号を区別する。\par

\noindent\textbf{原典照合：} 全5候補、表の列順、否定、条件、規格番号、法令時点を23原典画像で照合。\par

\paragraph{問われる核心} 安全の許容基準の社会・技術依存性と、設計から情報提供までの対策順位を判定する。

\paragraph{必要知識} 安全、許容リスク、社会的価値、設計優先。

\paragraph{考え方・正答を選ぶ理由（査読待ち初稿）} 不適切な4を選ぶ。許容リスクの条件依存性と、基準が永久に一定という主張は両立しない。国別の安全文化の断定を解法の必須前提にはしない。

\Needspace{4\baselineskip}\paragraph{選択肢の検討（査読待ち初稿）}\mbox{}\par

\noindent\textbf{1：国別一般化要確認} 日本は絶対安全、国際規格は許容リスク。。絶対安全と許容リスクの概念差は学べるが、日本全体を絶対安全志向とする国別の全称は一次資料で確定していない。\par

\noindent\textbf{2：国別一般化要確認} 日本は注意訓練重視、国際規格は技術向上重視。。人間だけに安全を依存させない設計は重要である。ただし日本・国際の一律の対比は検証範囲外である。\par

\noindent\textbf{3：一般要求として適切} 一般的安全要求は新規の機械にも適用。。一般要求から個別機械へ安全側面を導入する。全ての個別仕様まで一律になるという意味ではない。\par

\noindent\textbf{4：不適切} 許容基準は技術が進歩しても変わらない。。許容性は社会の価値と所与の条件に依存し、新たな知見・技術や価値観の変化で見直され得る。固定不変ではない。\par

\noindent\textbf{5：順位として適切} 設計、安全防護、使用上の情報の順。。危険を設計で減らす検討を先に行い、防護と情報で補完する。注意書きだけで設計検討を省略しない。\par

\paragraph{再利用する考え方} 危険源、曝露、損害、評価基準、対策、残留リスクを分ける。全称・必須・予算だけ・優先順位なしという限定を検討し、出題当時と現在の規定を混同しない。

\paragraph{共通理論} \hyperref[theory:safety:acceptability]{第\ref*{theory:safety:acceptability}節}, \hyperref[theory:safety:hierarchy]{第\ref*{theory:safety:hierarchy}節}

\par 参考：\cite{saf_guide51,saf_ra2006}。

\paragraph{関連問題} \hyperref[q:h29-reg-apt-12]{H29-REG-APT-12}

\par 難易度 2（校準初稿）、優先度未評価。

\subsection{平成29年度 問12}
\label{q:h29-reg-apt-12}

\noindent\textbf{原典：} sources/H29/aptitude.pdf、PDF物理11頁。

\textbf{公式正答：} 3（single）。

\par\textbf{検討未確定：} 肢2の「アセスメントを実施することによってリスクは軽減」を、措置の実施を含まない評価だけの効果と読むと不正確。Guide51全文未取得のため規格条文全体との照合も保留。

\paragraph{問題の要旨} 安全の条件依存性、残留リスク、対策の階層と人間の誤りを整理する。

\Needspace{4\baselineskip}\paragraph{確認した条件・記述（独自要約）}\mbox{}\par

\noindent\textbf{判断方向：} 各肢の理由に記録。組合せ・個数問題では記述の真偽と回答番号を区別する。\par

\noindent\textbf{原典照合：} 全5候補、表の列順、否定、条件、規格番号、法令時点を23原典画像で照合。\par

\paragraph{問われる核心} 安全の条件依存性、残留リスク、対策の階層と人間の誤りを整理する。

\paragraph{必要知識} 安全、残留リスク、反復、人間の誤り、社会条件。

\paragraph{考え方・正答を選ぶ理由（査読待ち初稿）} 最も不適切な3を選ぶ。許容基準と残留リスクを区別し、設備を変える措置と評価・記録も分ける。2は低減活動全体の略記と解釈するが、調査だけで必ず低減されるという全称命題にはしない。

\Needspace{4\baselineskip}\paragraph{選択肢の検討（査読待ち初稿）}\mbox{}\par

\noindent\textbf{1：適切（条件内）} 安全は許容されないリスクがないこと。。安全をゼロリスクと同一にせず、受容可能性の基準を含む概念として使う。\par

\noindent\textbf{2：適切（条件内）} 調査で低減し残留リスクを確認・文書化。。評価そのものが危険を物理的に消すわけではない。低減措置の実施と確認を伴う活動として読む。\par

\noindent\textbf{3：不適切} 許容基準は時代や社会で変わらない。。許容性の前提が社会・技術の変化で変わるため不適切である。\par

\noindent\textbf{4：適切（条件内）} 本質的設計、防護、使用情報の順。。情報だけで設計・防護の検討を置き換えない優先順位である。\par

\noindent\textbf{5：適切（条件内）} 人は間違えるため設備側を安全化。。注意を安全成立の唯一の条件にせず、防護と使用者の教育を組み合わせる。\par

\paragraph{再利用する考え方} 危険源、曝露、損害、評価基準、対策、残留リスクを分ける。全称・必須・予算だけ・優先順位なしという限定を検討し、出題当時と現在の規定を混同しない。

\paragraph{共通理論} \hyperref[theory:safety:acceptability]{第\ref*{theory:safety:acceptability}節}, \hyperref[theory:safety:cycle]{第\ref*{theory:safety:cycle}節}, \hyperref[theory:safety:hierarchy]{第\ref*{theory:safety:hierarchy}節}

\par 参考：\cite{saf_ra2006,saf_guide51}。

\paragraph{関連問題} \hyperref[q:h23-reg-apt-05]{H23-REG-APT-05}

\par 難易度 2（校準初稿）、優先度未評価。

\section{特定・見積り分析・評価・対応と反復}

\subsection{令和2年度 問7}
\label{q:r02-reg-apt-07}

\noindent\textbf{原典：} sources/R02/aptitude.pdf、PDF物理7頁。

\textbf{公式正答：} 4（single）。

\par\textbf{検討未確定：} 原典図・全候補は2頁の画像で確認済み。Guide51原規格の全文取得不能のため規格条項の直接照合は未完。

\paragraph{問題の要旨} 対策後にも見積りと評価を反復し、残留リスクの妥当性を確認して文書化する。

\Needspace{4\baselineskip}\paragraph{確認した条件・記述（独自要約）}\mbox{}\par

\noindent\textbf{判断方向：} 各肢の理由に記録。組合せ・個数問題では記述の真偽と回答番号を区別する。\par

\noindent\textbf{原典照合：} 全5候補、表の列順、否定、条件、規格番号、法令時点を23原典画像で照合。\par

\paragraph{問われる核心} 対策後にも見積りと評価を反復し、残留リスクの妥当性を確認して文書化する。

\paragraph{必要知識} 見積りと評価、反復、対策後の残留、妥当性確認と文書化。

\paragraph{考え方・正答を選ぶ理由（査読待ち初稿）} 図のリスク低減の下は、低減後のリスクを見積り、評価し、残留リスクが許容可能か判定する流れである。いいえなら低減へ戻り、はいなら妥当性確認及び文書化を経て完了する。4を選ぶ。上側の使用者・危険源の同定と、下側の再評価を区別する。

\Needspace{4\baselineskip}\paragraph{選択肢の検討（査読待ち初稿）}\mbox{}\par

\noindent\textbf{1：不適合} 見積り・評価・発生リスク・妥当性確認文書化。ウが違う。対策後に残る危険を評価するので残留リスクである。\par

\noindent\textbf{2：不適合} 同定・評価・発生リスク・合理性確認記録化。ア、ウ、エが違う。対策後の見積り・残留性・妥当性確認という流れを満たさない。\par

\noindent\textbf{3：不適合} 見積り・検証・残留リスク・妥当性確認記録化。イとエが違う。先に大きさを見積り、基準に照らして評価し、最後に妥当性確認と文書化を行う。\par

\noindent\textbf{4：図の流れに適合・規格本文保留} 見積り・評価・残留リスク・妥当性確認文書化。ア見積り、イ評価、ウ残留リスク、エ妥当性確認及び文書化で、原典図の反復と完了の位置に合う。\par

\noindent\textbf{5：不適合} 同定・検証・発生リスク・合理性確認文書化。ア、イ、ウ、エが違う。危険源同定や検証という語を反復中の見積り・評価へ置き換えない。\par

\paragraph{再利用する考え方} 危険源、曝露、損害、評価基準、対策、残留リスクを分ける。全称・必須・予算だけ・優先順位なしという限定を検討し、出題当時と現在の規定を混同しない。

\paragraph{共通理論} \hyperref[theory:safety:cycle]{第\ref*{theory:safety:cycle}節}

\par 参考：\cite{saf_ra2006,saf_guide51}。

\paragraph{関連問題} \hyperref[q:r05-reg-apt-08]{R05-REG-APT-08}

\par 難易度 2（校準初稿）、優先度未評価。

\subsection{令和5年度 問8}
\label{q:r05-reg-apt-08}

\noindent\textbf{原典：} sources/R05/aptitude.pdf、PDF物理8頁。

\textbf{公式正答：} 2（single）。

\par\textbf{検討未確定：} 原典図と全候補は確認済み。ISO公式公開説明は特定分析評価対応等の体系を確認できたが、JIS Q31000:2019本文全文への図・条項の直接照合は未完。

\paragraph{問題の要旨} 組織リスクの特定、分析、評価と、評価後の対応の位置を区別する。

\Needspace{4\baselineskip}\paragraph{確認した条件・記述（独自要約）}\mbox{}\par

\noindent\textbf{判断方向：} 各肢の理由に記録。組合せ・個数問題では記述の真偽と回答番号を区別する。\par

\noindent\textbf{原典照合：} 全5候補、表の列順、否定、条件、規格番号、法令時点を23原典画像で照合。\par

\paragraph{問われる核心} 組織リスクの特定、分析、評価と、評価後の対応の位置を区別する。

\paragraph{必要知識} 特定分析評価の囲み、対応、協議、監視と記録、2018ISO/2019JIS。

\paragraph{考え方・正答を選ぶ理由（査読待ち初稿）} 2を選ぶ。目的・状況と判断基準を定め、何が起こり得るかを特定し、その原因・結果・大きさを分析し、基準に対して評価する。その後、対応を選択し実施する。図の外周のコミュニケーション・協議と監視・レビュー、下の記録・報告は継続する活動である。

\Needspace{4\baselineskip}\paragraph{選択肢の検討（査読待ち初稿）}\mbox{}\par

\noindent\textbf{1：不適合} 分析・評価・対応・管理。アで特定を欠き、対応をアセスメント内へ入れ、エに管理を置く点が違う。\par

\noindent\textbf{2：原典図に適合・規格全文保留} 特定・分析・評価・対応。ア特定、イ分析、ウ評価、エ対応で図の順序とアセスメントの囲みに合う。\par

\noindent\textbf{3：不適合} 特定・評価・対応・管理。分析が抜け、対応を評価と同じ囲み内へ入れる。\par

\noindent\textbf{4：不適合} 分析・特定・評価・対応。事象を特定する前に分析する順序が違う。\par

\noindent\textbf{5：不適合} 分析・評価・特定・管理。特定を最後に置き、対応を欠く。\par

\paragraph{再利用する考え方} 危険源、曝露、損害、評価基準、対策、残留リスクを分ける。全称・必須・予算だけ・優先順位なしという限定を検討し、出題当時と現在の規定を混同しない。

\paragraph{共通理論} \hyperref[theory:safety:cycle]{第\ref*{theory:safety:cycle}節}, \hyperref[theory:safety:definition]{第\ref*{theory:safety:definition}節}, \hyperref[theory:safety:standards]{第\ref*{theory:safety:standards}節}

\par 参考：\cite{saf_iso31000}。

\paragraph{関連問題} \hyperref[q:r02-reg-apt-07]{R02-REG-APT-07}

\par 難易度 2（校準初稿）、優先度未評価。

\section{目的への不確かさと結果・起こりやすさ}

\subsection{平成28年度 問8}
\label{q:h28-reg-apt-08}

\noindent\textbf{原典：} sources/H28/aptitude.pdf、PDF物理8頁。

\textbf{公式正答：} 4（single）。

\par\textbf{検討未確定：} 原典画像の前文はJIS Q31010:2012。JSA公開表紙はIEC/ISO31010:2009、Risk assessment techniquesでありISO31000の同一翻訳との前文に疑義。原典を黙って訂正しない。

\paragraph{問題の要旨} 目的への不確かさの影響という広い定義と、結果・確率による表現の両立を判定する。

\Needspace{4\baselineskip}\paragraph{確認した条件・記述（独自要約）}\mbox{}\par

\noindent\textbf{判断方向：} 各肢の理由に記録。組合せ・個数問題では記述の真偽と回答番号を区別する。\par

\noindent\textbf{原典照合：} 全5候補、表の列順、否定、条件、規格番号、法令時点を23原典画像で照合。\par

\paragraph{問われる核心} 目的への不確かさの影響という広い定義と、結果・確率による表現の両立を判定する。

\paragraph{必要知識} 目的、正負の逸脱、知識不足、結果と確率、規格番号と対応版。

\paragraph{考え方・正答を選ぶ理由（査読待ち初稿）} 最も不適切な4を選ぶ。原典の注記4自体が結果と起こりやすさの組合せによる表現を述べる。前文のJIS Q31010:2012はJSA表紙ではIEC/ISO31010:2009に対応するリスクアセスメント技法であり、ISO31000と同一視する前文の番号に疑義がある。

\Needspace{4\baselineskip}\paragraph{選択肢の検討（査読待ち初稿）}\mbox{}\par

\noindent\textbf{1：適切（条件内）} 好ましい影響と好ましくない影響を含む。。期待からの逸脱は正負の両方向を含む。安全で負を扱うことと矛盾しない。\par

\noindent\textbf{2：適切（条件内）} 目的を明確にする。。何を達成したいかがリスクを認定する基準となる。\par

\noindent\textbf{3：適切（条件内）} 情報・理解の不足も不確かさの原因。。知識不足による不確かさと、事象のばらつきを区別できる。学習・調査で前者を減らせる場合がある。\par

\noindent\textbf{4：不適切} 旧定義の結果と確率の組合せは新定義と矛盾し使えない。。結果と起こりやすさは広い定義の具体的表現として引き続き使える。定義が広がることは具体モデルを全面否定することではない。\par

\noindent\textbf{5：適切（条件内）} 方針、価値、文化を踏まえ活用。。組織状況に合わせるが、法令や安全要求を独自方針だけで免除する意味ではない。\par

\paragraph{再利用する考え方} 危険源、曝露、損害、評価基準、対策、残留リスクを分ける。全称・必須・予算だけ・優先順位なしという限定を検討し、出題当時と現在の規定を混同しない。

\paragraph{共通理論} \hyperref[theory:safety:definition]{第\ref*{theory:safety:definition}節}, \hyperref[theory:safety:standards]{第\ref*{theory:safety:standards}節}

\par 参考：\cite{saf_iso31000,saf_jis31010}。

\par 難易度 2（校準初稿）、優先度未評価。

\section{リスク認識の心理・社会条件}

\subsection{平成23年度 問12}
\label{q:h23-reg-apt-12}

\noindent\textbf{原典：} sources/H23/aptitude.pdf、PDF物理13頁。

\textbf{公式正答：} 2（single）。

\par\textbf{検討未確定：} アの市民の理解不能、ウの特定個人の認知、キの損失枠の一律優位は一般的可能性・傾向に限定。全状況に成立することは一次資料で未確定。

\paragraph{問題の要旨} 同じ定量形式でも認知が一致するとは限らず、リスク認識の心理・社会条件を考慮する。

\Needspace{4\baselineskip}\paragraph{確認した条件・記述（独自要約）}\mbox{}\par

\noindent\textbf{判断方向：} 各肢の理由に記録。組合せ・個数問題では記述の真偽と回答番号を区別する。\par

\noindent\textbf{原典照合：} 全5候補、表の列順、否定、条件、規格番号、法令時点を23原典画像で照合。\par

\paragraph{問われる核心} 同じ定量形式でも認知が一致するとは限らず、リスク認識の心理・社会条件を考慮する。

\paragraph{必要知識} 客観推定と主観的認識、時間距離、自発性、利用可能性、提示枠。

\paragraph{考え方・正答を選ぶ理由（査読待ち初稿）} 初回はカのみ誤り、1個の2であり公式と一致する。アは専門化すると細部を理解しにくい状況が生じる可能性、イは社会文化条件、ウは時間距離、エは自発性、オは想起しやすさ、キは提示枠の効果を述べる。いずれも全員に必ず成立する法則とはしない。カは同じ積という形式から同じ入力・判断を導く点が誤りである。

\Needspace{4\baselineskip}\paragraph{選択肢の検討（査読待ち初稿）}\mbox{}\par

\noindent\textbf{1：不一致} 誤り0個。カの「認知バイアスがあっても同じ結論」が誤りなので0個ではない。\par

\noindent\textbf{2：初回・公式一致、全称に限定} 誤り1個。カのみを誤りとした初回判定である。積の形式が同じでも確率・被害の主観的評価が違えば結論は変わる。\par

\noindent\textbf{3：不一致} 誤り2個。ア〜オ・キを一般的な傾向や可能性として読むと、二つ目の明確な誤りはない。\par

\noindent\textbf{4：不一致} 誤り3個。三つの誤りを特定できない。直近の危険、本人の選択、思い出しやすさ、提示枠は認知に影響し得る。\par

\noindent\textbf{5：不一致} 誤り4個。四つの誤りを特定できない。社会的条件や価値を無視して客観値だけで認知を予測しない。\par

\paragraph{再利用する考え方} 危険源、曝露、損害、評価基準、対策、残留リスクを分ける。全称・必須・予算だけ・優先順位なしという限定を検討し、出題当時と現在の規定を混同しない。

\paragraph{共通理論} \hyperref[theory:safety:perception]{第\ref*{theory:safety:perception}節}, \hyperref[theory:safety:communication]{第\ref*{theory:safety:communication}節}

\par 参考：\cite{saf_fsc_psych,saf_fsc_glossary}。

\par 難易度 2（校準初稿）、優先度未評価。

\section{リスクの見積りモデルと数値化の限界}

\subsection{平成24年度 問11}
\label{q:h24-reg-apt-11}

\noindent\textbf{原典：} sources/H24/aptitude.pdf、PDF物理10頁。

\textbf{公式正答：} 3（single）。

\par\textbf{検討未確定：} 前文の1996年に国家規格JISを制定したとの規格史について当時のJSA一次資料未確認。TRとJISの別を推測訂正せず追加査読。

\paragraph{問題の要旨} リスクの定義と定量化手法を区別し、確率と損失の積を唯一の必須表現としない。

\Needspace{4\baselineskip}\paragraph{確認した条件・記述（独自要約）}\mbox{}\par

\noindent\textbf{判断方向：} 各肢の理由に記録。組合せ・個数問題では記述の真偽と回答番号を区別する。\par

\noindent\textbf{原典照合：} 全5候補、表の列順、否定、条件、規格番号、法令時点を23原典画像で照合。\par

\paragraph{問われる核心} リスクの定義と定量化手法を区別し、確率と損失の積を唯一の必須表現としない。

\paragraph{必要知識} 不確かさ、期待値、質的評価、特定分析評価、対応と効果。

\paragraph{考え方・正答を選ぶ理由（査読待ち初稿）} 最も不適切な3を選ぶ。2006年指針の見積りも、数値化必須ではなく行列・演算・グラフ等を例示する。定義、測定モデル、評価基準を別に扱う。前文の1996年のJIS制定という規格史は設問のまま保存し未確認とする。

\Needspace{4\baselineskip}\paragraph{選択肢の検討（査読待ち初稿）}\mbox{}\par

\noindent\textbf{1：適切（条件内）} 目的への不確かさの影響。。組織の目的と期待からの逸脱を結び付ける広い定義である。\par

\noindent\textbf{2：適切（条件内）} 安全では好ましくない結果の可能性。。安全を扱う際には危害という負の側面に注目する。広い組織リスクの正の側面を否定する意味ではない。\par

\noindent\textbf{3：不適切} 発生確率と重大さの積による定量表現が必須。。積による期待損失は一つのモデルである。質的分類、行列、加算スコア、リスクグラフもあり、必須ではない。\par

\noindent\textbf{4：適切（条件内）} 特定、分析、評価がアセスメント。。存在する事象を見つけ、構造と大きさを分析し、基準に照らして評価する。\par

\noindent\textbf{5：適切（条件内）} 対応は望む変化が実現することが重要。。措置を実施した事実と効果を区別し、実施後の状態を確認する。\par

\paragraph{再利用する考え方} 危険源、曝露、損害、評価基準、対策、残留リスクを分ける。全称・必須・予算だけ・優先順位なしという限定を検討し、出題当時と現在の規定を混同しない。

\paragraph{共通理論} \hyperref[theory:safety:quantification]{第\ref*{theory:safety:quantification}節}, \hyperref[theory:safety:definition]{第\ref*{theory:safety:definition}節}, \hyperref[theory:safety:cycle]{第\ref*{theory:safety:cycle}節}

\par 参考：\cite{saf_ra2006,saf_iso31000}。

\par 難易度 2（校準初稿）、優先度未評価。

\section{設計・防護・管理・情報の優先階層}

\subsection{平成30年度 問11}
\label{q:h30-reg-apt-11}

\noindent\textbf{原典：} sources/H30/aptitude.pdf、PDF物理11頁。

\textbf{公式正答：} 3（single）。

\par\textbf{検討未確定：} 肢2の化学物質義務は指定物質等に限定。肢4の重大危害の例外と肢5の「安衛法で規定」という法源の表現を補う必要がある。全肢の無限定妥当性とはしない。

\paragraph{問題の要旨} 職場の低減措置の順位と、一般の努力義務・対象化学物質の義務を分ける。

\Needspace{4\baselineskip}\paragraph{確認した条件・記述（独自要約）}\mbox{}\par

\noindent\textbf{判断方向：} 各肢の理由に記録。組合せ・個数問題では記述の真偽と回答番号を区別する。\par

\noindent\textbf{原典照合：} 全5候補、表の列順、否定、条件、規格番号、法令時点を23原典画像で照合。\par

\paragraph{問われる核心} 職場の低減措置の順位と、一般の努力義務・対象化学物質の義務を分ける。

\paragraph{必要知識} 指針10順位、57条の3対象範囲、著しい不均衡と重大危害、具体的実施時期。

\paragraph{考え方・正答を選ぶ理由（査読待ち初稿）} 最も不適切な3を選ぶ。保護具を管理的対策より先の恒常的代替策に置く順序が違う。2・4・5の範囲と法源の限定も残し、公式の唯一の候補番号と各文言の精度を別に管理する。

\Needspace{4\baselineskip}\paragraph{選択肢の検討（査読待ち初稿）}\mbox{}\par

\noindent\textbf{1：適切（条件内）} 特定、見積り、受入れ評価。。危害の程度と可能性を組み合わせる。発生確率を精密に数値化できることが常に必須ではない。\par

\noindent\textbf{2：適切（条件内）} 一般は努力、化学物質は調査義務。。2016年施行の化学物質調査義務は指定物質・含有率等の対象条件を伴う。全化学物質へ無条件に拡張しない。\par

\noindent\textbf{3：不適切} 設計、工学、防護具、管理の順。。2006年指針は設計・工学・管理・個人保護具の順である。第3と第4が逆。\par

\noindent\textbf{4：適切（条件内）} 低減費用が効果より著しく大きく不合理なら省略。。著しい不均衡は単に費用が便益を少し上回ることではない。死亡・重い後遺障害等で実施が著しく不合理といえない場合は措置が必要である。\par

\noindent\textbf{5：適切（条件内）} 建物、設備材料、手順等の変更で調査。。指針・関係規則に具体化される時期であり、列挙が安衛法本則そのものに書かれているとの意味には注意する。\par

\paragraph{再利用する考え方} 危険源、曝露、損害、評価基準、対策、残留リスクを分ける。全称・必須・予算だけ・優先順位なしという限定を検討し、出題当時と現在の規定を混同しない。

\paragraph{共通理論} \hyperref[theory:safety:hierarchy]{第\ref*{theory:safety:hierarchy}節}, \hyperref[theory:safety:law]{第\ref*{theory:safety:law}節}, \hyperref[theory:safety:alarp]{第\ref*{theory:safety:alarp}節}

\par 参考：\cite{saf_ra2006,saf_chem2016,saf_chemqa}。

\paragraph{関連問題} \hyperref[q:r06-reg-apt-13]{R06-REG-APT-13}

\par 難易度 2（校準初稿）、優先度未評価。

\subsection{令和元年度再試験 問12}
\label{q:r01-ret-apt-12}

\noindent\textbf{原典：} sources/R01\_retest/aptitude.pdf、PDF物理13頁。

\textbf{公式正答：} 2（single）。

\par\textbf{検討未確定：} Guide51全文取得不能のためア〜オの規格全条項への直接照合は未完。オを全ての水準に共通する唯一の許容条件へ拡張しない。

\paragraph{問題の要旨} 安全達成の反復、予見可能な誤使用、設計の優先順位とALARP条件を判定する。

\Needspace{4\baselineskip}\paragraph{確認した条件・記述（独自要約）}\mbox{}\par

\noindent\textbf{判断方向：} 各肢の理由に記録。組合せ・個数問題では記述の真偽と回答番号を区別する。\par

\noindent\textbf{原典照合：} 全5候補、表の列順、否定、条件、規格番号、法令時点を23原典画像で照合。\par

\paragraph{問われる核心} 安全達成の反復、予見可能な誤使用、設計の優先順位とALARP条件を判定する。

\paragraph{必要知識} 使用者群、予見可能な誤使用、残留リスク、階層、ALARP水準。

\paragraph{考え方・正答を選ぶ理由（査読待ち初稿）} 初回はエのみの1個、2で公式と一致する。設計、防護、使用情報には優先順位があり、まず設計検討を省略することは推奨されない。オはALARPを使う水準で、合理的に可能な追加低減を尽くした場合の条件として読む。

\Needspace{4\baselineskip}\paragraph{選択肢の検討（査読待ち初稿）}\mbox{}\par

\noindent\textbf{1：不一致} 不適切0個。エが本質安全設計の検討を省略できるとしており、0個ではない。\par

\noindent\textbf{2：初回・公式一致、規格本文照合保留} 不適切1個。エのみが明確に誤る。アは非ゼロリスクの安全、イは使用者と誤使用の同定、ウは反復と残留確認、オは適用領域内での合理的低減として読む。\par

\noindent\textbf{3：不一致} 不適切2個。他に二つ目の明確な誤りはない。ただしオの適用領域の限定を守る。\par

\noindent\textbf{4：不一致} 不適切3個。三つの誤りを特定できない。\par

\noindent\textbf{5：不一致} 不適切4個。四つの誤りを特定できない。\par

\paragraph{再利用する考え方} 危険源、曝露、損害、評価基準、対策、残留リスクを分ける。全称・必須・予算だけ・優先順位なしという限定を検討し、出題当時と現在の規定を混同しない。

\paragraph{共通理論} \hyperref[theory:safety:hierarchy]{第\ref*{theory:safety:hierarchy}節}, \hyperref[theory:safety:cycle]{第\ref*{theory:safety:cycle}節}, \hyperref[theory:safety:alarp]{第\ref*{theory:safety:alarp}節}

\par 参考：\cite{saf_ra2006,saf_guide51,saf_parliament2006}。

\paragraph{関連問題} \hyperref[q:r04-reg-apt-06]{R04-REG-APT-06}

\par 難易度 2（校準初稿）、優先度未評価。

\subsection{令和4年度 問6}
\label{q:r04-reg-apt-06}

\noindent\textbf{原典：} sources/R04/aptitude.pdf、PDF物理5頁。

\textbf{公式正答：} 5（single）。

\par\textbf{検討未確定：} Guide51全文の該当条項は未取得であり、規格本文への全5肢の直接照合は保留。公開一次資料で階層とALARP水準の原理を確認。

\paragraph{問題の要旨} ALARPの適用領域・著しい不均衡と、設計の優先順位を判定する。

\Needspace{4\baselineskip}\paragraph{確認した条件・記述（独自要約）}\mbox{}\par

\noindent\textbf{判断方向：} 各肢の理由に記録。組合せ・個数問題では記述の真偽と回答番号を区別する。\par

\noindent\textbf{原典照合：} 全5候補、表の列順、否定、条件、規格番号、法令時点を23原典画像で照合。\par

\paragraph{問われる核心} ALARPの適用領域・著しい不均衡と、設計の優先順位を判定する。

\paragraph{必要知識} 本質設計の優先、合理的可能性、3水準、反復と記録。

\paragraph{考え方・正答を選ぶ理由（査読待ち初稿）} 最も不適切な5を選ぶ。3・4はALARPを適用するリスク水準に限定して読む。許容不能な高リスクを予算だけで許容に変えず、残留が軽微な水準と条件付許容も区別する。

\Needspace{4\baselineskip}\paragraph{選択肢の検討（査読待ち初稿）}\mbox{}\par

\noindent\textbf{1：適切（条件内）} 残留リスクまで反復し確認・文書化。。評価と対策を反復し、最終的に残る危険の根拠を確認して記録する。\par

\noindent\textbf{2：適切（条件内）} 合理的に実行可能な限り低く。。略語は単に数学的絶対最小値を求める意味ではなく、合理的実行可能性を含む。\par

\noindent\textbf{3：適切（条件内）} 適用水準で合理的に可能な低減を行う。。条件付で許容する領域で更なる合理的低減を検討する。\par

\noindent\textbf{4：適切（条件内）} 追加低減が不能又は著しく不釣合いなら許容。。適用領域内で合理的に可能な措置が残っていないことを根拠とする。どの高リスクも費用だけで許容できる意味ではない。\par

\noindent\textbf{5：不適切} 設計、防護、情報の3方策に順位なし。。本質的設計を先に検討し、防護、情報へ続く順位がある。\par

\paragraph{再利用する考え方} 危険源、曝露、損害、評価基準、対策、残留リスクを分ける。全称・必須・予算だけ・優先順位なしという限定を検討し、出題当時と現在の規定を混同しない。

\paragraph{共通理論} \hyperref[theory:safety:hierarchy]{第\ref*{theory:safety:hierarchy}節}, \hyperref[theory:safety:alarp]{第\ref*{theory:safety:alarp}節}, \hyperref[theory:safety:cycle]{第\ref*{theory:safety:cycle}節}

\par 参考：\cite{saf_ra2006,saf_parliament2006,saf_guide51}。

\paragraph{関連問題} \hyperref[q:r01-ret-apt-12]{R01-RET-APT-12}, \hyperref[q:r07-reg-apt-14]{R07-REG-APT-14}

\par 難易度 2（校準初稿）、優先度未評価。

\section{リスク回避・低減・移転・保有と説明}

\subsection{平成24年度 問10}
\label{q:h24-reg-apt-10}

\noindent\textbf{原典：} sources/H24/aptitude.pdf、PDF物理9頁。

\textbf{公式正答：} 1（single）。

\par\textbf{検討未確定：} アの「受容すべき」を無条件の受容義務とする根拠は未確認。設計想定の限界と社会的受容を区別し公式1を保持。

\paragraph{問題の要旨} 設計外力を超える残留リスク、対応類型、未事故でも行う調査と説明を整理する。

\Needspace{4\baselineskip}\paragraph{確認した条件・記述（独自要約）}\mbox{}\par

\noindent\textbf{判断方向：} 各肢の理由に記録。組合せ・個数問題では記述の真偽と回答番号を区別する。\par

\noindent\textbf{原典照合：} 全5候補、表の列順、否定、条件、規格番号、法令時点を23原典画像で照合。\par

\paragraph{問われる核心} 設計外力を超える残留リスク、対応類型、未事故でも行う調査と説明を整理する。

\paragraph{必要知識} 設計条件、残留リスク、回避低減移転保有、当時の努力義務。

\paragraph{考え方・正答を選ぶ理由（査読待ち初稿）} 初回は○○○○、1で公式と一致する。対策設計には想定と限界を示す必要があり、重大さと目的に応じて対応を選ぶ。アの「受容すべき」は設計時に残る範囲を認識する趣旨として読み、住民が無条件で受忍しなければならないという規範へ拡張しない。

\Needspace{4\baselineskip}\paragraph{選択肢の検討（査読待ち初稿）}\mbox{}\par

\noindent\textbf{1：初回・公式一致、アの限定要確認} ア〜エ：○○○○。アは設計条件を超える事象を想定する、イは回避・低減・移転・保有、ウは予防的な調査の努力義務、エは認識への助言で、当時の設問意図では全て○。\par

\noindent\textbf{2：不一致} ×○○○。アを一律に誤りとはできない。ただし残留リスクの存在が無条件の社会的受容を正当化しない。\par

\noindent\textbf{3：不一致} ×○×○。アとウを誤りにする。災害がまだないことは危険がない証拠ではなく、当時の努力義務を否定できない。\par

\noindent\textbf{4：不一致} ○×○○。イの対応類型を否定するが、低減以外に回避や財務的移転・保有も選択肢となる。\par

\noindent\textbf{5：不一致} ○○○×。エを否定するが、情報提供と助言は必要である。説得して受け入れさせる一方通行だけを意味しない。\par

\paragraph{再利用する考え方} 危険源、曝露、損害、評価基準、対策、残留リスクを分ける。全称・必須・予算だけ・優先順位なしという限定を検討し、出題当時と現在の規定を混同しない。

\paragraph{共通理論} \hyperref[theory:safety:response]{第\ref*{theory:safety:response}節}, \hyperref[theory:safety:communication]{第\ref*{theory:safety:communication}節}, \hyperref[theory:safety:law]{第\ref*{theory:safety:law}節}

\par 参考：\cite{saf_ra2006,saf_fsc_glossary}。

\paragraph{関連問題} \hyperref[q:h26-reg-apt-10]{H26-REG-APT-10}

\par 難易度 2（校準初稿）、優先度未評価。

\subsection{平成26年度 問10}
\label{q:h26-reg-apt-10}

\noindent\textbf{原典：} sources/H26/aptitude.pdf、PDF物理10頁。

\textbf{公式正答：} 4（single）。

\paragraph{問題の要旨} 残留リスクの被害と不確かさを説明し、知らないことを説明不要の根拠にしない。

\Needspace{4\baselineskip}\paragraph{確認した条件・記述（独自要約）}\mbox{}\par

\noindent\textbf{判断方向：} 各肢の理由に記録。組合せ・個数問題では記述の真偽と回答番号を区別する。\par

\noindent\textbf{原典照合：} 全5候補、表の列順、否定、条件、規格番号、法令時点を23原典画像で照合。\par

\paragraph{問われる核心} 残留リスクの被害と不確かさを説明し、知らないことを説明不要の根拠にしない。

\paragraph{必要知識} 残留リスク、対応類型、予測限界と説明、2014時点努力義務。

\paragraph{考え方・正答を選ぶ理由（査読待ち初稿）} ○×○○の4を選ぶ。残る危険の発生条件と想定被害を伝え、不確かな範囲は不確かと明示する。未来を完全に予測できないことは、既知の被害や合理的に想定される結果を説明しなくてよいという意味ではない。

\Needspace{4\baselineskip}\paragraph{選択肢の検討（査読待ち初稿）}\mbox{}\par

\noindent\textbf{1：不一致} ア〜エ：○××○。アとエは適切だが、ウを×にする点が違う。\par

\noindent\textbf{2：不一致} ×○○×。ア・エを否定し、被害説明を免除するイを肯定する。\par

\noindent\textbf{3：不一致} ×○×○。適切なアを否定し、不適切なイを肯定する。\par

\noindent\textbf{4：適切な組合せ} ○×○○。アは対応の選択、イは説明不要として誤り、ウは当時の努力義務、エは助言で適切。\par

\noindent\textbf{5：不一致} ○×○×。エを×にする点が違う。市民の認識を考慮した情報交換は必要である。\par

\paragraph{再利用する考え方} 危険源、曝露、損害、評価基準、対策、残留リスクを分ける。全称・必須・予算だけ・優先順位なしという限定を検討し、出題当時と現在の規定を混同しない。

\paragraph{共通理論} \hyperref[theory:safety:response]{第\ref*{theory:safety:response}節}, \hyperref[theory:safety:communication]{第\ref*{theory:safety:communication}節}, \hyperref[theory:safety:law]{第\ref*{theory:safety:law}節}

\par 参考：\cite{saf_ra2006,saf_fsc_psych}。

\paragraph{関連問題} \hyperref[q:h24-reg-apt-10]{H24-REG-APT-10}

\par 難易度 2（校準初稿）、優先度未評価。

\section{食品リスク評価・管理・対話の役割}

\subsection{平成26年度 問11}
\label{q:h26-reg-apt-11}

\noindent\textbf{原典：} sources/H26/aptitude.pdf、PDF物理11頁。

\textbf{公式正答：} 1（single）。

\par\textbf{検討未確定：} 当時の全関係者文言の逐語照合は未実施。初回説明が用量評価と取り違えた事実をfollowupに保持し、改めて原典の役割逆転を説明。

\paragraph{問題の要旨} 食品リスクの科学的評価と管理・対話の役割を分ける。

\Needspace{4\baselineskip}\paragraph{確認した条件・記述（独自要約）}\mbox{}\par

\noindent\textbf{原典：} 問題文・否定・全五候補と組合せ表を保存原典27画像で確認。\par

\noindent\textbf{時点：} 平成26年度の本文表現を保持。後の制度改正を遡及しない。\par

\paragraph{問われる核心} 食品リスクの科学的評価と管理・対話の役割を分ける。

\paragraph{必要知識} 評価と管理の機能分担、リスクアナリシスの3要素、行政・事業者・消費者・専門家の相互作用。

\paragraph{考え方・正答を選ぶ理由（査読待ち初稿）} アはリスク評価を行政・事業者、管理を科学者中心とした役割逆転なので誤。科学的な健康影響の評価と、施策を検討・実施する管理を分ける。国・自治体の情報と対話の場、事業者の連携、消費者の参加、科学者の説明というイ〜オは正で1。対話は一方向の説得や安心の宣言だけではない。

\Needspace{4\baselineskip}\paragraph{選択肢の検討（査読待ち初稿）}\mbox{}\par

\noindent\textbf{1：条件内で適切} XOOOO。各文又は定義の判定を先に行う。全項目が一致する。\par

\noindent\textbf{2：不適切} OXXOX。各文又は定義の判定を先に行う。並びが判定 XOOOO と異なる。\par

\noindent\textbf{3：不適切} OOXOX。各文又は定義の判定を先に行う。並びが判定 XOOOO と異なる。\par

\noindent\textbf{4：不適切} OXOOX。各文又は定義の判定を先に行う。並びが判定 XOOOO と異なる。\par

\noindent\textbf{5：不適切} XOOXO。各文又は定義の判定を先に行う。並びが判定 XOOOO と異なる。\par

\paragraph{再利用する考え方} 対象・責任主体・時点・要件を分け、文の否定と「全て」「不要」の範囲を点検する。知識を当てはめる前に何を証明する資料かを確かめる。

\paragraph{共通理論} \hyperref[theory:safety:foodroles]{第\ref*{theory:safety:foodroles}節}

\par 参考：\cite{saf53_food_roles2007,saf_fsc_psych,saf_fsc_glossary}。

\par 難易度 3（校準初稿）、優先度未評価。

\section{GHSの分類・伝達と国内追加表示}

\subsection{平成24年度 問6}
\label{q:h24-reg-apt-06}

\noindent\textbf{原典：} sources/H24/aptitude.pdf、PDF物理5頁。

\textbf{公式正答：} 3（single）。

\par\textbf{検討未確定：} 当時のJISの全規格本文と全優先規則の版照合は未実施。H24公示と旧GHS案内で枠組みを確認し、細目はAstra対象。

\paragraph{問題の要旨} GHSは危険有害性の分類と伝達を調和する。

\Needspace{4\baselineskip}\paragraph{確認した条件・記述（独自要約）}\mbox{}\par

\noindent\textbf{原典：} 問題文・否定・全五候補と組合せ表を保存原典27画像で確認。\par

\noindent\textbf{時点：} 平成24年度の本文表現を保持。後の制度改正を遡及しない。\par

\paragraph{問われる核心} GHSは危険有害性の分類と伝達を調和する。

\paragraph{必要知識} 分類クラス・区分、ラベルとSDS、注意喚起語と絵表示、追加の法定情報。

\paragraph{考え方・正答を選ぶ理由（査読待ち初稿）} アは分類基準と表示の共通化、イは国内法による追加表示、ウは絵表示・注意喚起語、エは複数のラベル要素の優先規則を述べ、4文とも正で3。危険有害性の分類を、その物質のあらゆる使用条件でのリスク計算と同一視しない。

\Needspace{4\baselineskip}\paragraph{選択肢の検討（査読待ち初稿）}\mbox{}\par

\noindent\textbf{1：不適切} XOOO。各文又は定義の判定を先に行う。並びが判定 OOOO と異なる。\par

\noindent\textbf{2：不適切} OXOO。各文又は定義の判定を先に行う。並びが判定 OOOO と異なる。\par

\noindent\textbf{3：条件内で適切} OOOO。各文又は定義の判定を先に行う。全項目が一致する。\par

\noindent\textbf{4：不適切} OOXO。各文又は定義の判定を先に行う。並びが判定 OOOO と異なる。\par

\noindent\textbf{5：不適切} OOOX。各文又は定義の判定を先に行う。並びが判定 OOOO と異なる。\par

\paragraph{再利用する考え方} 対象・責任主体・時点・要件を分け、文の否定と「全て」「不要」の範囲を点検する。知識を当てはめる前に何を証明する資料かを確かめる。

\paragraph{共通理論} \hyperref[theory:safety:ghs]{第\ref*{theory:safety:ghs}節}

\par 参考：\cite{saf53_ghs_old,saf53_ghs_notice}。

\paragraph{関連問題} \hyperref[q:h26-reg-apt-08]{H26-REG-APT-08}

\par 難易度 3（校準初稿）、優先度未評価。

\subsection{平成26年度 問8}
\label{q:h26-reg-apt-08}

\noindent\textbf{原典：} sources/H26/aptitude.pdf、PDF物理8頁。

\textbf{公式正答：} 3（single）。

\par\textbf{検討未確定：} 当時のJIS規格全文の版照合は未実施。国内法の追加表示と更新の原則の範囲を確認した部分校正。

\paragraph{問題の要旨} GHSラベルに国内の法定事項を追加できる。

\Needspace{4\baselineskip}\paragraph{確認した条件・記述（独自要約）}\mbox{}\par

\noindent\textbf{原典：} 問題文・否定・全五候補と組合せ表を保存原典27画像で確認。\par

\noindent\textbf{時点：} 平成26年度の本文表現を保持。後の制度改正を遡及しない。\par

\paragraph{問われる核心} GHSラベルに国内の法定事項を追加できる。

\paragraph{必要知識} ラベル要素、供給者識別、新知見、JISと国内法の役割。

\paragraph{考え方・正答を選ぶ理由（査読待ち初稿）} アの物質名・供給者・危険有害性、ウの絵表示、エの新知見に応じた更新は正。イの国内法の追加表示を認めないという否定が誤なのでO X O Oの3。GHSは国内の別の表示義務を消すものではない。

\Needspace{4\baselineskip}\paragraph{選択肢の検討（査読待ち初稿）}\mbox{}\par

\noindent\textbf{1：不適切} OXXO。各文又は定義の判定を先に行う。並びが判定 OXOO と異なる。\par

\noindent\textbf{2：不適切} XOXO。各文又は定義の判定を先に行う。並びが判定 OXOO と異なる。\par

\noindent\textbf{3：条件内で適切} OXOO。各文又は定義の判定を先に行う。全項目が一致する。\par

\noindent\textbf{4：不適切} OOOX。各文又は定義の判定を先に行う。並びが判定 OXOO と異なる。\par

\noindent\textbf{5：不適切} OOOO。各文又は定義の判定を先に行う。並びが判定 OXOO と異なる。\par

\paragraph{再利用する考え方} 対象・責任主体・時点・要件を分け、文の否定と「全て」「不要」の範囲を点検する。知識を当てはめる前に何を証明する資料かを確かめる。

\paragraph{共通理論} \hyperref[theory:safety:ghs]{第\ref*{theory:safety:ghs}節}

\par 参考：\cite{saf53_ghs_old,saf53_ghs_notice}。

\paragraph{関連問題} \hyperref[q:h24-reg-apt-06]{H24-REG-APT-06}

\par 難易度 3（校準初稿）、優先度未評価。

\section{ヒューマンエラー発生要因の12分類}

\subsection{令和2年度 問8}
\label{q:r02-reg-apt-08}

\noindent\textbf{原典：} sources/R02/aptitude.pdf、PDF物理8頁。

\textbf{公式正答：} 1（single）。

\par\textbf{検討未確定：} 原典の9割以上という数値は統計母集団・集計法が未確認。疲労等を行為そのものと呼ぶ定義ではなく、原因分類として扱う。原典の該当しないの語義をAstra対象に保持。

\paragraph{問題の要旨} ヒューマンエラーを生む12の要因を結果からたどる。

\Needspace{4\baselineskip}\paragraph{確認した条件・記述（独自要約）}\mbox{}\par

\noindent\textbf{原典：} 問題文・否定・全五候補と組合せ表を保存原典27画像で確認。\par

\noindent\textbf{時点：} 令和2年度の本文表現を保持。後の制度改正を遡及しない。\par

\paragraph{問われる核心} ヒューマンエラーを生む12の要因を結果からたどる。

\paragraph{必要知識} エラーの行為と発生要因、認知・判断・操作、設備・作業手順・連絡・負荷・組織。

\paragraph{考え方・正答を選ぶ理由（査読待ち初稿）} 原典はエラーの発生要因を広く扱い、全12項目が該当する0の1。知識、注意、意思伝達、組織風土、近道行動、状況への反射、パニック、錯覚、心身機能、疲労、単調作業を別に調べる。誤操作という結果と、その前の条件を区別して対策を考える。

\Needspace{4\baselineskip}\paragraph{選択肢の検討（査読待ち初稿）}\mbox{}\par

\noindent\textbf{1：条件内で適切} 該当しない0項目。列挙は全てヒューマンエラー発生要因の分類例に対応するので、原典の用語法では0。\par

\noindent\textbf{2：不適切} 1項目。疲労や高齢化を単なる行為でないとして一つだけ除外すると、原典が使う要因分類と合わない。\par

\noindent\textbf{3：不適切} 2項目。心身機能低下・疲労等はエラーを生む条件として分類へ含まれる。行為の誤りとその条件は概念上分ける。\par

\noindent\textbf{4：不適切} 3項目。連絡不足や集団の条件も作業者を取り巻く要因で、個人のミスだけに狭めない。\par

\noindent\textbf{5：不適切} 4項目。12項目の分類例と対応し、4項目除外の根拠はない。\par

\paragraph{再利用する考え方} 対象・責任主体・時点・要件を分け、文の否定と「全て」「不要」の範囲を点検する。知識を当てはめる前に何を証明する資料かを確かめる。

\paragraph{共通理論} \hyperref[theory:safety:human]{第\ref*{theory:safety:human}節}

\par 参考：\cite{saf53_human12}。

\par 難易度 3（校準初稿）、優先度未評価。

\section{機械安全の方策とシステム境界}

\subsection{平成25年度 問5}
\label{q:h25-reg-apt-05}

\noindent\textbf{原典：} sources/H25/aptitude.pdf、PDF物理6頁。

\textbf{公式正答：} 3（single）。

\par\textbf{検討未確定：} 2007厚労省指針で方策を確認。原典が挙げる2003年ISO12100の全定義原文は未通読。初回説明のISO26000との取り違えをfollowupに記録し、正答一致だけでその説明を正当化しない。

\paragraph{問題の要旨} ISO12100のリスク低減は使用上の情報も含む。

\Needspace{4\baselineskip}\paragraph{確認した条件・記述（独自要約）}\mbox{}\par

\noindent\textbf{原典：} 問題文・否定・全五候補と組合せ表を保存原典27画像で確認。\par

\noindent\textbf{時点：} 平成25年度の本文表現を保持。後の制度改正を遡及しない。\par

\paragraph{問われる核心} ISO12100のリスク低減は使用上の情報も含む。

\paragraph{必要知識} 信頼性と安全性、リスク分析と評価、本質的安全→安全防護・付加保護→使用上の情報。

\paragraph{考え方・正答を選ぶ理由（査読待ち初稿）} 指定された機能を期間・条件内で果たす信頼性は正。危険源の同定、リスク見積り、評価を含むアセスメントも正。ウは本質的安全と防護の2方策へ限定して使用上の情報を落とすので誤。本質的安全はガードを追加する前に設計上の危険を減らすためエは正。O O X Oの3。

\Needspace{4\baselineskip}\paragraph{選択肢の検討（査読待ち初稿）}\mbox{}\par

\noindent\textbf{1：不適切} XOOO。各文又は定義の判定を先に行う。並びが判定 OOXO と異なる。\par

\noindent\textbf{2：不適切} OXOO。各文又は定義の判定を先に行う。並びが判定 OOXO と異なる。\par

\noindent\textbf{3：条件内で適切} OOXO。各文又は定義の判定を先に行う。全項目が一致する。\par

\noindent\textbf{4：不適切} OOOX。各文又は定義の判定を先に行う。並びが判定 OOXO と異なる。\par

\noindent\textbf{5：不適切} OOOO。各文又は定義の判定を先に行う。並びが判定 OOXO と異なる。\par

\paragraph{再利用する考え方} 対象・責任主体・時点・要件を分け、文の否定と「全て」「不要」の範囲を点検する。知識を当てはめる前に何を証明する資料かを確かめる。

\paragraph{共通理論} \hyperref[theory:safety:machine]{第\ref*{theory:safety:machine}節}

\par 参考：\cite{saf53_machine2007}。

\par 難易度 3（校準初稿）、優先度未評価。

\subsection{令和元年度 問6}
\label{q:r01-reg-apt-06}

\noindent\textbf{原典：} sources/R01/aptitude.pdf、PDF物理7頁。

\textbf{公式正答：} 1（single）。

\par\textbf{検討未確定：} 機能安全の電子機器・付加という原典表現は一般規格の全範囲を尽くさない。安全工学・レジリエンスの当時用語源を全条照合していないため部分校正。

\paragraph{問題の要旨} システム安全は環境・物・人を含み回復も考える。

\Needspace{4\baselineskip}\paragraph{確認した条件・記述（独自要約）}\mbox{}\par

\noindent\textbf{原典：} 問題文・否定・全五候補と組合せ表を保存原典27画像で確認。\par

\noindent\textbf{時点：} 令和元年度の本文表現を保持。後の制度改正を遡及しない。\par

\paragraph{問われる核心} システム安全は環境・物・人を含み回復も考える。

\paragraph{必要知識} システムの境界、機能安全と制御、本質的安全・残留リスク、レジリエンス。

\paragraph{考え方・正答を選ぶ理由（査読待ち初稿）} アはハードだけではないB。イは安全のための装置の機能が正しく働くA。ウは製造側が安全を図り、残留リスクを全て使用者へ渡すとはしないA。エは使用者・作業者両方の危害を扱うA。オは予防と回復力を扱うB。B A A A Bの1。原典の機能安全の対比は安全機能の目的で読み、制御装置を全て安全機能と呼ばない。

\Needspace{4\baselineskip}\paragraph{選択肢の検討（査読待ち初稿）}\mbox{}\par

\noindent\textbf{1：条件内で適切} BAAAB。各文又は定義の判定を先に行う。全項目が一致する。\par

\noindent\textbf{2：不適切} BBBBA。各文又は定義の判定を先に行う。並びが判定 BAAAB と異なる。\par

\noindent\textbf{3：不適切} AAABA。各文又は定義の判定を先に行う。並びが判定 BAAAB と異なる。\par

\noindent\textbf{4：不適切} ABAAB。各文又は定義の判定を先に行う。並びが判定 BAAAB と異なる。\par

\noindent\textbf{5：不適切} BAABA。各文又は定義の判定を先に行う。並びが判定 BAAAB と異なる。\par

\paragraph{再利用する考え方} 対象・責任主体・時点・要件を分け、文の否定と「全て」「不要」の範囲を点検する。知識を当てはめる前に何を証明する資料かを確かめる。

\paragraph{共通理論} \hyperref[theory:safety:machine]{第\ref*{theory:safety:machine}節}

\par 参考：\cite{saf53_machine2007,saf_guide51}。

\par 難易度 3（校準初稿）、優先度未評価。

\section{2021インフラ行動計画の七項目}

\subsection{令和5年度 問13}
\label{q:r05-reg-apt-13}

\noindent\textbf{原典：} sources/R05/aptitude.pdf、PDF物理13頁。

\textbf{公式正答：} 5（single）。

\paragraph{問題の要旨} 2021行動計画の七項目と人材施策の内容を分ける。

\Needspace{4\baselineskip}\paragraph{確認した条件・記述（独自要約）}\mbox{}\par

\noindent\textbf{原典：} 問題文・否定・全五候補と組合せ表を保存原典27画像で確認。\par

\noindent\textbf{時点：} 令和5年度の本文表現を保持。後の制度改正を遡及しない。\par

\paragraph{問われる核心} 2021行動計画の七項目と人材施策の内容を分ける。

\paragraph{必要知識} 計画の章・対象期間、七項目の見出しと内包する施策、維持管理のサイクル。

\paragraph{考え方・正答を選ぶ理由（査読待ち初稿）} 該当しない5。2021年6月18日の第2次計画Ⅵの七項目は、個別施設計画、点検診断・修繕更新、予算管理、体制、新技術、情報基盤、基準類である。前章の現状課題には法令等の項目もあるため、章を取り違えない。技術継承が施策として不要という答えではない。

\Needspace{4\baselineskip}\paragraph{選択肢の検討（査読待ち初稿）}\mbox{}\par

\noindent\textbf{1：条件内で適切} 予算管理。維持更新費用の見通しと財政支援等を扱う独立項目に含まれる。\par

\noindent\textbf{2：条件内で適切} 体制の構築。担い手・役割分担・支援等の体制を整える独立項目に含まれる。\par

\noindent\textbf{3：条件内で適切} 新技術の開発・導入。点検・維持更新の技術と適用の評価を扱う独立項目に含まれる。\par

\noindent\textbf{4：条件内で適切} 情報基盤の整備と活用。施設情報の蓄積・更新と利用を扱う独立項目に含まれる。\par

\noindent\textbf{5：不適切} 技術継承の取組。重要性の否定ではなく、原典が問う七つの独立見出しとして列挙されていない。体制・人材の中の内容との区別が必要。\par

\paragraph{再利用する考え方} 対象・責任主体・時点・要件を分け、文の否定と「全て」「不要」の範囲を点検する。知識を当てはめる前に何を証明する資料かを確かめる。

\paragraph{共通理論} \hyperref[theory:safety:maintenance2021]{第\ref*{theory:safety:maintenance2021}節}

\par 参考：\cite{saf53_infra2021}。

\par 難易度 2（校準初稿）、優先度未評価。

\section{重大製品事故の報告・公表と時点}

\subsection{平成29年度 問9}
\label{q:h29-reg-apt-09}

\noindent\textbf{原典：} sources/H29/aptitude.pdf、PDF物理9頁。

\textbf{公式正答：} 3（single）。

\par\textbf{検討未確定：} 2017年当時の全除外品目の逐条対照は未実施。9品目から後年の対象変更を区別する資料を確認した部分校正。

\paragraph{問題の要旨} 重大製品事故の調査を理由に報告を遅らせない。

\Needspace{4\baselineskip}\paragraph{確認した条件・記述（独自要約）}\mbox{}\par

\noindent\textbf{原典：} 問題文・否定・全五候補と組合せ表を保存原典27画像で確認。\par

\noindent\textbf{時点：} 平成29年度の本文表現を保持。後の制度改正を遡及しない。\par

\paragraph{問われる核心} 重大製品事故の調査を理由に報告を遅らせない。

\paragraph{必要知識} 重大製品事故の報告と公表、消費生活用の範囲、長期使用の点検、出題時の指定数。

\paragraph{考え方・正答を選ぶ理由（査読待ち初稿）} 不適切3。事故の原因と製品起因の判断を調べつつ、報告に必要な事実と未確定事項を区別して期限内に伝える。調査中は何も知らせなくてよいという制度ではない。原典の9品目は当時の指定として扱う。

\Needspace{4\baselineskip}\paragraph{選択肢の検討（査読待ち初稿）}\mbox{}\par

\noindent\textbf{1：条件内で適切} 任意収集から法定報告への制度変更。重大事故の情報を行政が把握して再発防止へ使うため法定報告制度が導入された。法全体の創設年と事故報告改正年を分ける。\par

\noindent\textbf{2：条件・時点を留保} 一般消費者用を基本とし他法の対象を除外。消費生活用の範囲と政令等による除外を確認する。他法で安全性が担保されるという一般説明だけから全除外を列挙したものとは扱わない。\par

\noindent\textbf{3：不適切} 分析が済めば期限を超えてもよい。不明点のある段階でも所定期限内の報告が必要。調査の完成を期限延長の根拠にしない。\par

\noindent\textbf{4：条件・時点を留保} 特定保守製品9品目。2017年当時の指定を述べる。2021年の対象見直し等を遡及せず、現在の品目数と区別する。\par

\noindent\textbf{5：条件内で適切} 設計標準使用期間・点検期間を設定。長期使用に伴う劣化に対して点検の情報と期間を整える。耐用の無条件保証や点検だけで全危険が消える意味ではない。\par

\paragraph{再利用する考え方} 対象・責任主体・時点・要件を分け、文の否定と「全て」「不要」の範囲を点検する。知識を当てはめる前に何を証明する資料かを確かめる。

\paragraph{共通理論} \hyperref[theory:safety:productaccident]{第\ref*{theory:safety:productaccident}節}

\par 参考：\cite{saf53_products2012,saf53_oldnine}。

\paragraph{関連問題} \hyperref[q:h30-reg-apt-10]{H30-REG-APT-10}

\par 難易度 3（校準初稿）、優先度未評価。

\subsection{平成30年度 問10}
\label{q:h30-reg-apt-10}

\noindent\textbf{原典：} sources/H30/aptitude.pdf、PDF物理10頁。

\textbf{公式正答：} 5（single）。

\par\textbf{検討未確定：} 原典の1週間を法定の全案件絶対期限としない。知った日算入は取得2023手引きで確認したが2018施行版の当該府令全文は未照合。2009の制度施行と法改正時点も原典表現と分けて査読。

\paragraph{問題の要旨} 事故報告10日と公表の運用を別の時計で読む。

\Needspace{4\baselineskip}\paragraph{確認した条件・記述（独自要約）}\mbox{}\par

\noindent\textbf{原典：} 問題文・否定・全五候補と組合せ表を保存原典27画像で確認。\par

\noindent\textbf{時点：} 平成30年度の本文表現を保持。後の制度改正を遡及しない。\par

\paragraph{問われる核心} 事故報告10日と公表の運用を別の時計で読む。

\paragraph{必要知識} 重大事故の類型、事故発生と知った日の差、事業者の報告と行政の公表、通常の安全性。

\paragraph{考え方・正答を選ぶ理由（査読待ち初稿）} アの死亡・火災・CO中毒・後遺障害・治療30日以上、イの製造輸入側の報告、ウの公表、エの通常の安全性は原典の趣旨で全正5。報告は知った日を含む10日、公表は報告受理後の原則1週間等で起点が違う。2013公表フローはガス石油機器や製品起因が疑われるものを直ちに公表し、他の段階を原則1週間とする。

\Needspace{4\baselineskip}\paragraph{選択肢の検討（査読待ち初稿）}\mbox{}\par

\noindent\textbf{1：不適切} XOOO。各文又は定義の判定を先に行う。並びが判定 OOOO と異なる。\par

\noindent\textbf{2：不適切} OXOO。各文又は定義の判定を先に行う。並びが判定 OOOO と異なる。\par

\noindent\textbf{3：不適切} OOXO。各文又は定義の判定を先に行う。並びが判定 OOOO と異なる。\par

\noindent\textbf{4：不適切} OOOX。各文又は定義の判定を先に行う。並びが判定 OOOO と異なる。\par

\noindent\textbf{5：条件内で適切} OOOO。各文又は定義の判定を先に行う。全項目が一致する。\par

\paragraph{再利用する考え方} 対象・責任主体・時点・要件を分け、文の否定と「全て」「不要」の範囲を点検する。知識を当てはめる前に何を証明する資料かを確かめる。

\paragraph{共通理論} \hyperref[theory:safety:productaccident]{第\ref*{theory:safety:productaccident}節}

\par 参考：\cite{saf53_products2012,saf53_productflow2013,saf53_products2023}。

\paragraph{関連問題} \hyperref[q:h29-reg-apt-09]{H29-REG-APT-09}

\par 難易度 3（校準初稿）、優先度未評価。

\section{品質不正の記録と組織・現場条件}

\subsection{平成28年度 問5}
\label{q:h28-reg-apt-05}

\noindent\textbf{原典：} sources/H28/aptitude.pdf、PDF物理5頁。

\textbf{公式正答：} 4（single）。

\par\textbf{検討未確定：} 2015中間報告の公開概要と本文該当枠組みを確認。全41頁の事案・技術措置を通読した確認ではない。

\paragraph{問題の要旨} くい工事の記録不正と構造安全を個別に検証する。

\Needspace{4\baselineskip}\paragraph{確認した条件・記述（独自要約）}\mbox{}\par

\noindent\textbf{原典：} 問題文・否定・全五候補と組合せ表を保存原典27画像で確認。\par

\noindent\textbf{時点：} 平成28年度の本文表現を保持。後の制度改正を遡及しない。\par

\paragraph{問われる核心} くい工事の記録不正と構造安全を個別に検証する。

\paragraph{必要知識} 施工記録の真正性、現地条件との照合、元請の統括と下請の専門判断、異常時の手順。

\paragraph{考え方・正答を選ぶ理由（査読待ち初稿）} アの記録不正を直ちに構造不安全と断定しない点は、放置してよいという意味ではない。迅速に実際の支持層到達・安全性を確認し、不正の是正も行う。イの組織風土とウの元請・下請の責任体制は正。エの現場条件を無視した設計図優先、オのハード改善から人的ルール不要を導く主張は誤。O O O X Xの4。

\Needspace{4\baselineskip}\paragraph{選択肢の検討（査読待ち初稿）}\mbox{}\par

\noindent\textbf{1：不適切} OOOOX。各文又は定義の判定を先に行う。並びが判定 OOOXX と異なる。\par

\noindent\textbf{2：不適切} OXOXX。各文又は定義の判定を先に行う。並びが判定 OOOXX と異なる。\par

\noindent\textbf{3：不適切} XOOXO。各文又は定義の判定を先に行う。並びが判定 OOOXX と異なる。\par

\noindent\textbf{4：条件内で適切} OOOXX。各文又は定義の判定を先に行う。全項目が一致する。\par

\noindent\textbf{5：不適切} XOXOO。各文又は定義の判定を先に行う。並びが判定 OOOXX と異なる。\par

\paragraph{再利用する考え方} 対象・責任主体・時点・要件を分け、文の否定と「全て」「不要」の範囲を点検する。知識を当てはめる前に何を証明する資料かを確かめる。

\paragraph{共通理論} \hyperref[theory:safety:quality]{第\ref*{theory:safety:quality}節}

\par 参考：\cite{saf53_pile2015}。

\paragraph{関連問題} \hyperref[q:r01-reg-apt-07]{R01-REG-APT-07}

\par 難易度 3（校準初稿）、優先度未評価。

\subsection{令和元年度 問7}
\label{q:r01-reg-apt-07}

\noindent\textbf{原典：} sources/R01/aptitude.pdf、PDF物理8頁。

\textbf{公式正答：} 5（single）。

\par\textbf{検討未確定：} 2020官庁白書の2017年以降の複数事案記述で背景を確認。2019設問の全事案別調査報告の照合は未実施。

\paragraph{問題の要旨} 品質不正を個人だけへ帰して組織要因を消さない。

\Needspace{4\baselineskip}\paragraph{確認した条件・記述（独自要約）}\mbox{}\par

\noindent\textbf{原典：} 問題文・否定・全五候補と組合せ表を保存原典27画像で確認。\par

\noindent\textbf{時点：} 令和元年度の本文表現を保持。後の制度改正を遡及しない。\par

\paragraph{問われる核心} 品質不正を個人だけへ帰して組織要因を消さない。

\paragraph{必要知識} 記録の真正性、経営の優先順位、独立した検査・内部監査、部門間の責任と情報。

\paragraph{考え方・正答を選ぶ理由（査読待ち初稿）} アの他社の教訓は参考にならない、イの単発に限定、エの倫理活動があるから個人だけの問題という3文が誤。トップの判断・圧力と責任分担・他部門への配慮は組織条件に関わるのでウ・オは正。X X O X Oの5。個人の行為の責任を消すのではなく、行為を生む仕組みも改善する。

\Needspace{4\baselineskip}\paragraph{選択肢の検討（査読待ち初稿）}\mbox{}\par

\noindent\textbf{1：不適切} XOOXO。各文又は定義の判定を先に行う。並びが判定 XXOXO と異なる。\par

\noindent\textbf{2：不適切} XXXXX。各文又は定義の判定を先に行う。並びが判定 XXOXO と異なる。\par

\noindent\textbf{3：不適切} XOOOO。各文又は定義の判定を先に行う。並びが判定 XXOXO と異なる。\par

\noindent\textbf{4：不適切} OOOOO。各文又は定義の判定を先に行う。並びが判定 XXOXO と異なる。\par

\noindent\textbf{5：条件内で適切} XXOXO。各文又は定義の判定を先に行う。全項目が一致する。\par

\paragraph{再利用する考え方} 対象・責任主体・時点・要件を分け、文の否定と「全て」「不要」の範囲を点検する。知識を当てはめる前に何を証明する資料かを確かめる。

\paragraph{共通理論} \hyperref[theory:safety:quality]{第\ref*{theory:safety:quality}節}

\par 参考：\cite{saf53_quality2020,saf53_pile2015}。

\paragraph{関連問題} \hyperref[q:h28-reg-apt-05]{H28-REG-APT-05}

\par 難易度 3（校準初稿）、優先度未評価。

\section{リコールと安全規制・責任の区別}

\subsection{平成25年度 問8}
\label{q:h25-reg-apt-08}

\noindent\textbf{原典：} sources/H25/aptitude.pdf、PDF物理9頁。

\textbf{公式正答：} 2（single）。

\par\textbf{検討未確定：} 2010リコール手引き直接取得は未成功。H25の全対象追跡時の社告省略の細条件は未確定で、後年手引きだけで当時の全文を確認済みにしない。

\paragraph{問題の要旨} リコールと安全規制・PL責任を分ける。

\Needspace{4\baselineskip}\paragraph{確認した条件・記述（独自要約）}\mbox{}\par

\noindent\textbf{原典：} 問題文・否定・全五候補と組合せ表を保存原典27画像で確認。\par

\noindent\textbf{時点：} 平成25年度の本文表現を保持。後の制度改正を遡及しない。\par

\paragraph{問われる核心} リコールと安全規制・PL責任を分ける。

\paragraph{必要知識} 回収・修理等の是正、対象製品の追跡、行政報告、危害防止と民事責任。

\paragraph{考え方・正答を選ぶ理由（査読待ち初稿）} 誤り2は医薬品の品質管理制度から回収の不存在を導いた飛躍。回収は危害防止の是正で、事故がまだ起きていなくても対象となり得る。設計・劣化・表示を検討し、行政報告、購入者への連絡、社告、PL賠償を別の要件として扱う。

\Needspace{4\baselineskip}\paragraph{選択肢の検討（査読待ち初稿）}\mbox{}\par

\noindent\textbf{1：条件内で適切} 経年劣化による事故も対象。設計時の欠陥に限らず、使用で劣化し重大な危害を生む製品も是正の対象になり得る。\par

\noindent\textbf{2：不適切} 医薬品は品質管理法があるためリコール対象外。医薬品にも回収制度がある。他の安全規制があることは回収が不要である根拠ではない。消安法の対象外と、医薬品制度での回収を区別する。\par

\noindent\textbf{3：条件・時点を留保} 全対象を追跡回収できる場合の報告・社告。行政への報告と広く知らせる社告の要否は別。全対象が確実に捕捉され危害を防げる場合を前提に、当時のリコール手引きと照合する。\par

\noindent\textbf{4：条件内で適切} 表示誤りもリコール対象となり得る。危害を防ぐ注意事項や成分等の表示誤りは、実体の故障がなくても是正の対象となり得る。\par

\noindent\textbf{5：条件内で適切} 回収が直ちにPL責任を意味しない。回収判断と損害賠償の成立は別。PL責任では欠陥・損害・因果関係と対象等を検討する。\par

\paragraph{再利用する考え方} 対象・責任主体・時点・要件を分け、文の否定と「全て」「不要」の範囲を点検する。知識を当てはめる前に何を証明する資料かを確かめる。

\paragraph{共通理論} \hyperref[theory:safety:recall]{第\ref*{theory:safety:recall}節}

\par 参考：\cite{saf53_recall2010,saf53_medicine_recall,pl_statute,pl_qa}。

\paragraph{関連問題} \hyperref[q:h27-reg-apt-13]{H27-REG-APT-13}

\par 難易度 3（校準初稿）、優先度未評価。

\subsection{平成27年度 問13}
\label{q:h27-reg-apt-13}

\noindent\textbf{原典：} sources/H27/aptitude.pdf、PDF物理13頁。

\textbf{公式正答：} 2（single）。

\par\textbf{検討未確定：} 2010リコール手引き直接取得は未成功。H25の全対象追跡時の社告省略の細条件は未確定で、後年手引きだけで当時の全文を確認済みにしない。

\paragraph{問題の要旨} リコールと安全規制・PL責任を分ける。

\Needspace{4\baselineskip}\paragraph{確認した条件・記述（独自要約）}\mbox{}\par

\noindent\textbf{原典：} 問題文・否定・全五候補と組合せ表を保存原典27画像で確認。\par

\noindent\textbf{時点：} 平成27年度の本文表現を保持。後の制度改正を遡及しない。\par

\paragraph{問われる核心} リコールと安全規制・PL責任を分ける。

\paragraph{必要知識} 回収・修理等の是正、対象製品の追跡、行政報告、危害防止と民事責任。

\paragraph{考え方・正答を選ぶ理由（査読待ち初稿）} 誤り2は医薬品の品質管理制度から回収の不存在を導いた飛躍。回収は危害防止の是正で、事故がまだ起きていなくても対象となり得る。設計・劣化・表示を検討し、行政報告、購入者への連絡、社告、PL賠償を別の要件として扱う。

\Needspace{4\baselineskip}\paragraph{選択肢の検討（査読待ち初稿）}\mbox{}\par

\noindent\textbf{1：条件内で適切} 経年劣化による事故も対象。設計時の欠陥に限らず、使用で劣化し重大な危害を生む製品も是正の対象になり得る。\par

\noindent\textbf{2：不適切} 医薬品は品質管理法があるためリコール対象外。医薬品にも回収制度がある。他の安全規制があることは回収が不要である根拠ではない。消安法の対象外と、医薬品制度での回収を区別する。\par

\noindent\textbf{3：条件内で適切} 全対象が回収済みでも行政への報告が必要。回収したことと所管機関への報告義務は別。報告の目的は対象把握や再発防止にもある。\par

\noindent\textbf{4：条件内で適切} 表示誤りもリコール対象となり得る。危害を防ぐ注意事項や成分等の表示誤りは、実体の故障がなくても是正の対象となり得る。\par

\noindent\textbf{5：条件内で適切} 回収が直ちにPL責任を意味しない。回収判断と損害賠償の成立は別。PL責任では欠陥・損害・因果関係と対象等を検討する。\par

\paragraph{再利用する考え方} 対象・責任主体・時点・要件を分け、文の否定と「全て」「不要」の範囲を点検する。知識を当てはめる前に何を証明する資料かを確かめる。

\paragraph{共通理論} \hyperref[theory:safety:recall]{第\ref*{theory:safety:recall}節}

\par 参考：\cite{saf53_recall2010,saf53_medicine_recall,pl_statute,pl_qa}。

\paragraph{関連問題} \hyperref[q:h25-reg-apt-08]{H25-REG-APT-08}

\par 難易度 3（校準初稿）、優先度未評価。

\section{変更時のリスク調査と軽微例外}

\subsection{令和元年度 問11}
\label{q:r01-reg-apt-11}

\noindent\textbf{原典：} sources/R01/aptitude.pdf、PDF物理12頁。

\textbf{公式正答：} 2（single）。

\paragraph{問題の要旨} 変更時のリスク調査と明らかな軽微作業の例外。

\Needspace{4\baselineskip}\paragraph{確認した条件・記述（独自要約）}\mbox{}\par

\noindent\textbf{原典：} 問題文・否定・全五候補と組合せ表を保存原典27画像で確認。\par

\noindent\textbf{時点：} 令和元年度の本文表現を保持。後の制度改正を遡及しない。\par

\paragraph{問われる核心} 変更時のリスク調査と明らかな軽微作業の例外。

\paragraph{必要知識} 危険源同定、変更管理、予見可能性、重篤度と可能性、ALARPと著しい不均衡。

\paragraph{考え方・正答を選ぶ理由（査読待ち初稿）} アは建設物・材料・手順等が変わる時の再評価で正。イは明らかに軽微な負傷等しか予想されない歩行まで絶対除外不可としたため誤。ウの疲労の付加影響とエのALARPは正、2。明らかな軽微の例外を、不明な危険を調べなくてよい理由へ使わない。ALARPは安ければ十分という原則ではない。

\Needspace{4\baselineskip}\paragraph{選択肢の検討（査読待ち初稿）}\mbox{}\par

\noindent\textbf{1：不適切} OXXO。各文又は定義の判定を先に行う。並びが判定 OXOO と異なる。\par

\noindent\textbf{2：条件内で適切} OXOO。各文又は定義の判定を先に行う。全項目が一致する。\par

\noindent\textbf{3：不適切} OOXX。各文又は定義の判定を先に行う。並びが判定 OXOO と異なる。\par

\noindent\textbf{4：不適切} OOOX。各文又は定義の判定を先に行う。並びが判定 OXOO と異なる。\par

\noindent\textbf{5：不適切} XXOO。各文又は定義の判定を先に行う。並びが判定 OXOO と異なる。\par

\paragraph{再利用する考え方} 対象・責任主体・時点・要件を分け、文の否定と「全て」「不要」の範囲を点検する。知識を当てはめる前に何を証明する資料かを確かめる。

\paragraph{共通理論} \hyperref[theory:safety:workplace]{第\ref*{theory:safety:workplace}節}

\par 参考：\cite{saf53_workplace2006,saf_parliament2006}。

\par 難易度 2（校準初稿）、優先度未評価。

\section{職場リスク調査の対象・実施・統括・措置}

\subsection{平成23年度 問6}
\label{q:h23-reg-apt-06}

\noindent\textbf{原典：} sources/H23/aptitude.pdf、PDF物理7頁。

\textbf{公式正答：} 4（single）。

\paragraph{問題の要旨} 2006年のリスク調査の努力義務、実施時期、統括と現場参加を区別する。

\Needspace{4\baselineskip}\paragraph{確認した条件・記述（独自要約）}\mbox{}\par

\noindent\textbf{判断方向：} 各肢の理由に記録。組合せ・個数問題では記述の真偽と回答番号を区別する。\par

\noindent\textbf{原典照合：} 全5候補、表の列順、否定、条件、規格番号、法令時点を23原典画像で照合。\par

\paragraph{問われる核心} 2006年のリスク調査の努力義務、実施時期、統括と現場参加を区別する。

\paragraph{必要知識} 2006年法28条の2、指針4・5、現場参加と統括。

\paragraph{考え方・正答を選ぶ理由（査読待ち初稿）} 最も不適切な4を選ぶ。現場情報を集めることは必要であるが、統括管理まで現場だけに任せる根拠にならない。2・3は当時の指針を読み、5は適切な実施と措置を前提とした効果として理解する。

\Needspace{4\baselineskip}\paragraph{選択肢の検討（査読待ち初稿）}\mbox{}\par

\noindent\textbf{1：適切（条件内）} 特定、見積り、優先度、措置の決定。。危険を具体化し対策を選ぶ一連の手順である。評価書作成だけで設備の危険が消えるわけではない。\par

\noindent\textbf{2：適切（条件内）} 製造・建設等は調査と措置が努力義務。。2006年の一般制度の対象業種等の限定を伴う。全ての個別危害防止規定が努力義務だけになる意味ではない。\par

\noindent\textbf{3：適切（条件内）} 設備・材料・手順等の変更時に調査。。2006年指針5の時点で調査し、開始前の低減を求める。計画段階の検討は効果的である。\par

\noindent\textbf{4：不適切} 統括管理を現場の責任で行う。。指針4は事業の実施を統括管理する者に管理させる。労働者・職長・専門家の参加と統括責任の移譲は違う。\par

\noindent\textbf{5：適切（条件内）} 調査が説明責任と経営リスク低減へつながる。。適切な調査と措置・記録は危険の説明と経営上の管理を支える。実施という事実だけで製造物責任が免除されるものではない。\par

\paragraph{再利用する考え方} 危険源、曝露、損害、評価基準、対策、残留リスクを分ける。全称・必須・予算だけ・優先順位なしという限定を検討し、出題当時と現在の規定を混同しない。

\paragraph{共通理論} \hyperref[theory:safety:governance]{第\ref*{theory:safety:governance}節}, \hyperref[theory:safety:law]{第\ref*{theory:safety:law}節}, \hyperref[theory:safety:cycle]{第\ref*{theory:safety:cycle}節}

\par 参考：\cite{saf_ra2006,saf_ra_effects}。

\par 難易度 2（校準初稿）、優先度未評価。

\subsection{令和3年度 問15}
\label{q:r03-reg-apt-15}

\noindent\textbf{原典：} sources/R03/aptitude.pdf、PDF物理15頁。

\textbf{公式正答：} 3（single）。

\par\textbf{検討未確定：} オは公式では×を含むが、専門家の分析によって危険度が明確になるという一般命題は必ずしも偽でない。資料の効果列挙との不一致を科学的否定へ作り替えない。前文の化学物質義務にも指定対象の限定が必要。

\paragraph{問題の要旨} 厚労省の導入効果の表現と、残留リスクの規則・職場参加の効果を照合する。

\Needspace{4\baselineskip}\paragraph{確認した条件・記述（独自要約）}\mbox{}\par

\noindent\textbf{判断方向：} 各肢の理由に記録。組合せ・個数問題では記述の真偽と回答番号を区別する。\par

\noindent\textbf{原典照合：} 全5候補、表の列順、否定、条件、規格番号、法令時点を23原典画像で照合。\par

\paragraph{問われる核心} 厚労省の導入効果の表現と、残留リスクの規則・職場参加の効果を照合する。

\paragraph{必要知識} 残留リスクの管理規則、認識共有、全員参加の感受性と専門分析の別。

\paragraph{考え方・正答を選ぶ理由（査読待ち初稿）} 資料の記載を問う趣旨では○○○××、3である。厚労省はエに守るべき決め事の理由、オに全員参加による危険への感受性を挙げる。しかし専門家の分析で危険の大きさが明確になること自体はあり得る。引用リストとの不一致と命題の科学的偽を同一視せず、オの文言の疑義を残す。

\Needspace{4\baselineskip}\paragraph{選択肢の検討（査読待ち初稿）}\mbox{}\par

\noindent\textbf{1：記載と不一致} ア〜オ：○○○○○。原典のエ・オは厚労省の効果列挙と異なる。\par

\noindent\textbf{2：記載と不一致} ○○○○×。エは「守るべき決め事の理由」を「リスク発生要因の理由」に置換している。\par

\noindent\textbf{3：初回・公式一致、オの一般命題疑義} ○○○××。アはリスク明確化、イは共有、ウは優先度。エとオは資料の列挙との不一致として×。ただしオの一般命題まで否定できるかは保留。\par

\noindent\textbf{4：ウを誤否定} ○○×××。ウの合理的な優先順位の効果を否定してしまう。\par

\noindent\textbf{5：確認済み効果を誤否定} ×××××。ア・イ・ウという確認できた導入効果まで否定する。\par

\paragraph{再利用する考え方} 危険源、曝露、損害、評価基準、対策、残留リスクを分ける。全称・必須・予算だけ・優先順位なしという限定を検討し、出題当時と現在の規定を混同しない。

\paragraph{共通理論} \hyperref[theory:safety:participation]{第\ref*{theory:safety:participation}節}, \hyperref[theory:safety:governance]{第\ref*{theory:safety:governance}節}

\par 参考：\cite{saf_ra_effects,saf_ra2006,saf_chem2016}。

\par 難易度 3（校準初稿）、優先度未評価。

\subsection{令和6年度 問13}
\label{q:r06-reg-apt-13}

\noindent\textbf{原典：} sources/R06/aptitude.pdf、PDF物理13頁。

\textbf{公式正答：} 5（single）。

\paragraph{問題の要旨} 2006年指針の調査対象、疲労の付加影響、優先措置と重大危害への暫定措置を確認する。

\Needspace{4\baselineskip}\paragraph{確認した条件・記述（独自要約）}\mbox{}\par

\noindent\textbf{判断方向：} 各肢の理由に記録。組合せ・個数問題では記述の真偽と回答番号を区別する。\par

\noindent\textbf{原典照合：} 全5候補、表の列順、否定、条件、規格番号、法令時点を23原典画像で照合。\par

\paragraph{問われる核心} 2006年指針の調査対象、疲労の付加影響、優先措置と重大危害への暫定措置を確認する。

\paragraph{必要知識} 指針6調査選定、8疲労、10順位と暫定、軽微の限定。

\paragraph{考え方・正答を選ぶ理由（査読待ち初稿）} ○○○○の5を選ぶ。イの除外は単に以前軽傷だけだったことを意味せず、現時点の十分な検討で明らかに軽微な危害しか予想されないという条件である。ウは疲労や単調作業等で発生可能性が増す点を考慮する。エの費用の例外も重大危害を放置する許可ではなく、根本対策に時間を要するときは暫定措置を直ちに行う。

\Needspace{4\baselineskip}\paragraph{選択肢の検討（査読待ち初稿）}\mbox{}\par

\noindent\textbf{1：不一致} ア〜エ：×○○○。アの特定、見積り、優先度、措置、記録の定義を否定する点が違う。\par

\noindent\textbf{2：不一致} ○×○○。イは明らかに軽微しか予想されない場合の限定で除外できるため、×ではない。平坦という名称だけで全通路を除外しない。\par

\noindent\textbf{3：不一致} ○○×○。ウの疲労等の影響考慮は指針8にあるため×ではない。\par

\noindent\textbf{4：不一致} ○○○×。エの高順位対策、著しい不合理性の例外と、重大リスクの暫定措置の即時実施は指針10にあり×ではない。\par

\noindent\textbf{5：適切な組合せ} ○○○○。アからエまで2006年指針の定義と条件に対応する。\par

\paragraph{再利用する考え方} 危険源、曝露、損害、評価基準、対策、残留リスクを分ける。全称・必須・予算だけ・優先順位なしという限定を検討し、出題当時と現在の規定を混同しない。

\paragraph{共通理論} \hyperref[theory:safety:law]{第\ref*{theory:safety:law}節}, \hyperref[theory:safety:hierarchy]{第\ref*{theory:safety:hierarchy}節}, \hyperref[theory:safety:cycle]{第\ref*{theory:safety:cycle}節}

\par 参考：\cite{saf_ra2006}。

\paragraph{関連問題} \hyperref[q:h30-reg-apt-11]{H30-REG-APT-11}

\par 難易度 2（校準初稿）、優先度未評価。

\section{職場安全の経験比・報告・4S・化学情報}

\subsection{令和3年度 問12}
\label{q:r03-reg-apt-12}

\noindent\textbf{原典：} sources/R03/aptitude.pdf、PDF物理12頁。

\textbf{公式正答：} 3（single）。

\par\textbf{検討未確定：} カは全化学物質・全使用者に無条件の法的調査義務があるようにも読め、指定物質等の対象範囲を補う必要がある。キも一般健診報告は事業場規模・健診種類に依存。

\paragraph{問題の要旨} 労働安全の目的、事故統計の限界、報告保護、4S、化学情報と健康診断を区別する。

\Needspace{4\baselineskip}\paragraph{確認した条件・記述（独自要約）}\mbox{}\par

\noindent\textbf{判断方向：} 各肢の理由に記録。組合せ・個数問題では記述の真偽と回答番号を区別する。\par

\noindent\textbf{原典照合：} 全5候補、表の列順、否定、条件、規格番号、法令時点を23原典画像で照合。\par

\paragraph{問われる核心} 労働安全の目的、事故統計の限界、報告保護、4S、化学情報と健康診断を区別する。

\paragraph{必要知識} 目的と責任体制、状態/行動、経験比と因果、報告保護、4S、SDS、健診。

\paragraph{考え方・正答を選ぶ理由（査読待ち初稿）} 初回はア・イ・エ・カ・キの5個、3で公式と一致する。ウの1:29:300は経験的集計の教訓で、軽傷原因の除去が重災害ゼロの必要十分条件になるわけではない。オの4Sは整理・整頓・清掃・清潔である。カの調査義務は対象物質等、キの一般健診報告は事業場規模等に限定される。

\Needspace{4\baselineskip}\paragraph{選択肢の検討（査読待ち初稿）}\mbox{}\par

\noindent\textbf{1：不足} 適切3個。ア・イ・エ・カ・キを適切とする初回の5個に足りない。\par

\noindent\textbf{2：不足} 適切4個。同じ5個に足りない。\par

\noindent\textbf{3：初回・公式一致、法の対象条件要確認} 適切5個。アは目的達成の手段、イは状態と行動、エは報告者を責めない、カは対象物質のSDS活用、キは健診・報告の対象条件内で適切。ウの軽傷要因除去だけで重災害をなくせるという保証と、オの4Sへしつけを入れる点が誤り。\par

\noindent\textbf{4：過剰} 適切6個。ウ・オのいずれかを肯定することになり不適合。\par

\noindent\textbf{5：過剰} 適切7個。ウもオも肯定することになり不適合。\par

\paragraph{再利用する考え方} 危険源、曝露、損害、評価基準、対策、残留リスクを分ける。全称・必須・予算だけ・優先順位なしという限定を検討し、出題当時と現在の規定を混同しない。

\paragraph{共通理論} \hyperref[theory:safety:incident]{第\ref*{theory:safety:incident}節}, \hyperref[theory:safety:housekeeping]{第\ref*{theory:safety:housekeeping}節}, \hyperref[theory:safety:law]{第\ref*{theory:safety:law}節}

\par 参考：\cite{saf_welfare,saf_4s,saf_chemqa,saf_healthreport}。

\paragraph{関連問題} \hyperref[q:r06-reg-apt-07]{R06-REG-APT-07}

\par 難易度 3（校準初稿）、優先度未評価。

\subsection{令和6年度 問7}
\label{q:r06-reg-apt-07}

\noindent\textbf{原典：} sources/R06/aptitude.pdf、PDF物理7頁。

\textbf{公式正答：} 4（single）。

\par\textbf{検討未確定：} イの「同じ人間330件」「背後の数千」の経験集計の原著・母集団条件は未照合。厚労省資料で1:29:300と非難しない活動を確認したが、全個人の普遍則とはしない。

\paragraph{問題の要旨} 職場の状態・行動、経験的事故比、報告保護、4Sと化学情報を判断する。

\Needspace{4\baselineskip}\paragraph{確認した条件・記述（独自要約）}\mbox{}\par

\noindent\textbf{判断方向：} 各肢の理由に記録。組合せ・個数問題では記述の真偽と回答番号を区別する。\par

\noindent\textbf{原典照合：} 全5候補、表の列順、否定、条件、規格番号、法令時点を23原典画像で照合。\par

\paragraph{問われる核心} 職場の状態・行動、経験的事故比、報告保護、4Sと化学情報を判断する。

\paragraph{必要知識} 状態行動、経験的比、情報提供者の保護、4S、SDSの性質と取扱い。

\paragraph{考え方・正答を選ぶ理由（査読待ち初稿）} 初回は○○○×○、4で公式と一致する。エの4Sは清掃であり洗浄ではない。イの1:29:300は安全活動の教訓としての経験比で、各個人やあらゆる業種の件数が必ず一致する確率法則ではない。SDSの一般的機能と法的交付の指定対象は別に整理する。

\Needspace{4\baselineskip}\paragraph{選択肢の検討（査読待ち初稿）}\mbox{}\par

\noindent\textbf{1：不一致} ア〜オ：○○○××。SDSの一般的役割であるオを否定する点が違う。\par

\noindent\textbf{2：不一致} ○○××○。報告者を責めないウを否定する点が違う。\par

\noindent\textbf{3：不一致} ×××○○。適切なア・ウ・オを否定し、誤るエを肯定する。\par

\noindent\textbf{4：初回・公式一致、経験比条件保留} ○○○×○。アは状態と行動、イは伝統的経験比、ウは報告保護、エは洗浄を入れた誤り、オはSDSの情報伝達機能で○○○×○。\par

\noindent\textbf{5：不一致} ×××○×。状態・行動と報告保護を否定し4Sの誤った列挙を肯定する。\par

\paragraph{再利用する考え方} 危険源、曝露、損害、評価基準、対策、残留リスクを分ける。全称・必須・予算だけ・優先順位なしという限定を検討し、出題当時と現在の規定を混同しない。

\paragraph{共通理論} \hyperref[theory:safety:incident]{第\ref*{theory:safety:incident}節}, \hyperref[theory:safety:housekeeping]{第\ref*{theory:safety:housekeeping}節}, \hyperref[theory:safety:law]{第\ref*{theory:safety:law}節}

\par 参考：\cite{saf_welfare,saf_4s,saf_chemqa}。

\paragraph{関連問題} \hyperref[q:r03-reg-apt-12]{R03-REG-APT-12}

\par 難易度 2（校準初稿）、優先度未評価。
